
\documentclass[UTF8,openany,zihao=-4]{ctexbook}
\usepackage[a4paper,top=1.45cm,bottom=1.45cm,left=1.35cm,right=1.35cm,headsep=0.25cm,footskip=0.65cm]{geometry}
\usepackage{amsmath,amssymb}
\usepackage{enumitem}
\usepackage{booktabs,longtable,array,tabularx,multicol}
\usepackage{xcolor}
\usepackage{hyperref}
\usepackage{fancyhdr}
\usepackage{lastpage}
\usepackage{titlesec}
\hypersetup{colorlinks=true,linkcolor=black,urlcolor=black,citecolor=black}
\setlength{\parindent}{0em}
\setlength{\parskip}{0.18em}
\linespread{1.02}
\raggedbottom
\pagestyle{fancy}
\setlength{\headheight}{15pt}
\fancyhf{}
\fancyhead[L]{操作系统原理高密度理解讲义}
\fancyhead[R]{\leftmark}
\fancyfoot[C]{\thepage/\pageref{LastPage}}
\renewcommand{\headrulewidth}{0.25pt}
\ctexset{
  chapter = {format = \huge\bfseries},
  section = {format = \Large\bfseries},
  subsection = {format = \large\bfseries}
}
\titlespacing*{\chapter}{0pt}{-0.6em}{0.8em}
\titlespacing*{\section}{0pt}{0.55em}{0.22em}
\titlespacing*{\subsection}{0pt}{0.35em}{0.12em}
\newlist{points}{enumerate}{2}
\setlist[points,1]{label=\arabic*. ,leftmargin=2.0em,itemsep=0.06em,topsep=0.06em,parsep=0em,partopsep=0em}
\setlist[points,2]{label=(\arabic*) ,leftmargin=1.8em,itemsep=0.02em,topsep=0.02em,parsep=0em,partopsep=0em}
\setlist[itemize]{itemsep=0.04em,topsep=0.04em,parsep=0em,partopsep=0em}
\setlist[enumerate]{itemsep=0.04em,topsep=0.04em,parsep=0em,partopsep=0em}
\newcommand{\concept}[1]{\subsection{#1}}
\newcommand{\key}[1]{\textbf{#1}}
\newcommand{\term}[1]{\textbf{#1}}
\newcommand{\code}[1]{\texttt{#1}}
\newcommand{\zhushi}[2]{\par\noindent{\kaishu\small #1：#2}\par}
\newcommand{\plain}[1]{\zhushi{大白话}{#1}}
\newcommand{\exam}[1]{\zhushi{可能怎么考}{#1}}
\newcolumntype{Y}{>{\raggedright\arraybackslash}X}
\title{\bfseries 操作系统原理高密度理解讲义\\\large 按 PPT 顺序：知识点 + 大白话 + 可能怎么考}
\author{基于课程 PPT、往年试卷与个人 os.docx 辅助记忆整理}
\date{\today}

\begin{document}
\maketitle
\tableofcontents


\chapter*{使用说明}
\addcontentsline{toc}{chapter}{使用说明}
\begin{points}
  \item 本讲义按 PPT 顺序组织，尽量覆盖从绪论到设备管理的完整考试知识链。
  \item 每个知识点先用普通字体写“需要背和会用的知识”，再用楷体写“\textbf{大白话}”和“\textbf{可能怎么考}”。
  \item 背诵时先背普通字体；做题前看楷体部分，建立题目识别直觉。
  \item 版式改为高密度文本，不使用花哨色块；重点靠编号、粗体和楷体区分。
\end{points}
\chapter{操作系统绪论}

\section{操作系统的概念与定位}
\concept{计算机系统的层次}
\begin{points}
  \item 计算机系统由\term{硬件}和\term{软件}组成。硬件包括 CPU、内存、I/O 设备等；软件包括系统软件和应用软件。
  \item 操作系统是配置在裸机上的\key{第一层软件}，位于硬件与其他软件之间，是其他软件运行的基础。
  \item 编译器、数据库管理系统、编辑程序也属于系统软件，但它们不是操作系统本身。
  \item 从层次上看：硬件/裸机 $\rightarrow$ 操作系统 $\rightarrow$ 系统程序 $\rightarrow$ 应用程序 $\rightarrow$ 用户。
\end{points}
\plain{硬件不会直接服务应用，OS 是硬件上第一层“扩展机”，把裸机包装成可管理、可调用的机器。}
\exam{常考“谁是第一层软件”“编译器/数据库是不是 OS”“OS 位于哪两层之间”。}

\concept{为什么需要操作系统}
\begin{points}
  \item \term{从用户视角看}：操作系统是用户和计算机硬件之间的接口，为用户提供方便、安全、高效的使用环境。
  \item \term{从系统视角看}：操作系统是计算机系统资源的管理者，负责分配、回收、保护和调度资源。
  \item 操作系统通过三种思想扩展计算机能力：\key{资源复用、资源虚化、资源抽象}。
  \item 复用解决“资源数量不够”：CPU 轮流给多个进程用是时间复用；内存分块给多个进程用是空间复用。
  \item 虚化解决“一个物理资源变成多个逻辑资源”：一个 CPU 虚拟成多个虚拟处理机，一台打印机通过 SPOOLing 虚拟成多台逻辑打印机。
  \item 抽象解决“硬件太复杂”：磁盘扇区被抽象成文件，CPU 执行流被抽象成进程，物理设备被抽象成逻辑设备名。
\end{points}
\plain{没有 OS，用户要直接面对 CPU、内存、磁盘、设备；OS 的价值就是管资源、挡危险、把复杂硬件变简单。}
\exam{常考从用户视角/系统视角回答，或解释资源复用、虚化、抽象。}


\concept{操作系统定义}
\begin{points}
  \item 操作系统没有唯一的精确定义，但可从两个角度理解。
  \item \term{控制程序}：控制程序执行，防止错误和不当使用，协调应用与硬件。
  \item \term{资源管理器}：管理处理机、存储器、I/O 设备、文件、磁盘空间等软硬件资源。
  \item 简洁表述：操作系统是位于裸机和应用程序之间的系统软件，对下管理控制计算机软硬件资源，对上提供方便、安全、高效的接口和运行环境。
\end{points}
\plain{定义题不要追求一句神句，抓住“系统软件、资源管理者、用户接口/扩展机”三个关键词就稳。}
\exam{常考简答“什么是操作系统”，按“位置-对象-作用-接口”写。}


\section{操作系统的形成与发展}
\concept{发展阶段总览}
\begin{points}
  \item 操作系统发展顺序：\key{手工操作 $\rightarrow$ 批处理系统 $\rightarrow$ 多道批处理系统 $\rightarrow$ 分时系统 $\rightarrow$ 实时系统 $\rightarrow$ 通用操作系统}。
  \item 每个阶段都是在解决上一阶段的主要矛盾：人机矛盾、CPU 与 I/O 速度矛盾、交互需求、实时响应需求。
\end{points}
\plain{发展史就是不断解决矛盾：人机慢、CPU 等 I/O、用户要交互、任务有截止期限。}
\exam{常考按阶段排序，或问某阶段解决了什么问题。}


\concept{手工操作阶段}
\begin{points}
  \item 用户独占机器，人工装入程序、启动运行、取走结果。
  \item 主要矛盾：人太慢，CPU 太快，I/O 太慢，导致资源利用率很低。
  \item 没有真正的操作系统，只是人手工控制机器。
\end{points}
\plain{人手工装程序，CPU 大量等人和 I/O，所以效率极低。}
\exam{常考缺点：用户独占资源、CPU 利用率低、没有自动作业控制。}


\concept{批处理系统}
\begin{points}
  \item 把多个作业组织成一批，由监督程序 Monitor 自动控制，一个接一个执行。
  \item 优点：减少人工干预，提高吞吐量。
  \item 缺点：缺少交互性，用户提交作业后不能直接干预运行过程。
  \item 关键词：\key{自动、成批、无交互、吞吐量、周转时间}。
\end{points}
\plain{把一堆作业交给监督程序自动跑，减少人工干预，但用户不能边运行边改。}
\exam{常考关键词“成批、自动、无交互、吞吐量”。}


\concept{多道批处理系统}
\begin{points}
  \item 多个作业同时进入内存，交替运行；当一个作业等待 I/O 时，CPU 可执行另一个作业。
  \item 核心目标：解决 CPU 等待 I/O 的浪费，提高 CPU 和 I/O 设备利用率。
  \item 代价：需要进程调度、内存分配、资源保护、中断处理、上下文切换等机制，系统开销更大。
  \item 直觉例子：作业 A 等磁盘读入时，作业 B 用 CPU 计算；磁盘和 CPU 可以并行工作。
\end{points}
\plain{多个作业进内存，一个等 I/O，另一个用 CPU，让 CPU 和设备都别闲着。}
\exam{常考为什么能提高效率、为什么需要中断/调度/存储保护。}


\concept{分时系统}
\begin{points}
  \item 多个用户通过终端同时使用计算机，系统把 CPU 时间划分为时间片，轮流响应各用户请求。
  \item 追求目标：较快响应用户，使每个用户感觉自己独占机器。
  \item 特征：\key{多路性、交互性、独立性、及时性}。
  \item 注意：分时系统的及时性是“用户命令较快响应”，不是实时系统那种严格截止期限。
  \item 若时间片为 $Q$，用户数为 $N$，粗略响应周期可理解为 $T=Q\times N$。用户越多，每个用户轮到 CPU 的间隔越长。
\end{points}
\plain{把 CPU 切成很多时间片轮流给用户，目标是让每个终端都感觉机器在响应自己。}
\exam{常考分时四特征，尤其“及时性”不等于硬实时。}


\concept{实时系统}
\begin{points}
  \item 实时系统强调在规定时间内响应外部事件。不是单纯“速度快”，而是响应必须\key{及时、可预测、可靠}。
  \item 硬实时：错过截止时间会造成严重失败，例如飞控、安全控制。
  \item 软实时：偶尔超时可以容忍，例如多媒体播放中偶尔卡顿。
  \item 关键词：\key{及时性、可靠性、截止时间、硬实时、软实时}。
\end{points}
\plain{实时不是“跑得快”而是“截止时间前必须响应”。}
\exam{常考硬实时/软实时区别、实时系统追求可靠性和可预测性。}


\section{操作系统的基本特征}
\concept{并发性}
\begin{points}
  \item 并发指若干事件在\key{同一时间间隔内}发生。
  \item 并行指若干事件在\key{同一时刻}发生。
  \item 单 CPU 上多个进程可以并发，但不能并行；多 CPU/多核上多个进程可以真正并行。
  \item 课件常用判断：并行一定是并发，但并发不一定是并行；并行是并发的特例。
\end{points}
\plain{核心是“看起来同时推进”。单 CPU 上靠时间片交替，多 CPU 上才可能真的同一时刻运行。}
\exam{常考并发与并行区别、单处理机下最多几个进程运行、并发导致为什么需要同步互斥。}


\concept{共享性}
\begin{points}
  \item 共享指系统资源可供多个并发执行的进程共同使用。
  \item \term{互斥共享}：一段时间内只允许一个进程访问，例如打印机、临界区。
  \item \term{同时共享}：宏观上多个进程可同时访问，例如只读文件、共享代码段。
  \item 并发与共享互为条件：没有并发，就谈不上资源竞争；没有共享管理，并发程序也难以正确运行。
\end{points}
\plain{多个进程都想用同一批资源，所以必须规定能不能一起用、谁先用、用完怎么释放。}
\exam{常考互斥共享/同时共享举例，或者问并发和共享为什么是 OS 最基本特征。}


\concept{虚拟性}
\begin{points}
  \item 虚拟性指通过某种技术把一个物理实体变成若干逻辑上的对应物。
  \item 典型例子：虚拟处理机、虚拟存储器、虚拟设备。
  \item 虚拟不是“凭空变多”，而是通过复用、映射、换入换出等机制，让用户看到更方便的逻辑资源。
\end{points}
\plain{不是资源真的变多，而是 OS 用映射、换入换出、排队等办法，让用户感觉自己有一份独立资源。}
\exam{常考虚拟处理机、虚拟存储器、虚拟设备的例子，以及虚拟性和资源复用的关系。}


\concept{异步性/不确定性}
\begin{points}
  \item 异步性指并发环境下，程序以不可预知的速度向前推进。
  \item 一个进程何时被调度、何时暂停、I/O 等多久、何时继续执行，不完全由程序自己决定。
  \item 异步性带来问题：若没有同步互斥机制，可能出现数据不一致、结果不可再现。
\end{points}
\plain{进程什么时候跑、什么时候停，不完全由程序自己决定，而是被调度、中断、I/O 共同影响。}
\exam{常考并发程序为什么不可再现、为什么会有与时间有关的错误。}


\section{操作系统的功能与接口}
\concept{功能模块}
\begin{points}
  \item \term{处理机管理/进程管理}：进程控制、进程同步、进程通信、调度。
  \item \term{存储器管理}：内存分配与回收、地址映射、保护共享、存储扩充。
  \item \term{设备管理}：I/O 控制、设备分配、缓冲、设备驱动、设备独立性。
  \item \term{文件管理}：文件存储、目录、文件共享、文件保护、外存空间管理。
  \item \term{作业管理/用户接口}：作业控制、命令接口、系统调用接口、图形接口。
\end{points}
\plain{五大管理就是 OS 的主线：CPU、内存、设备、文件、用户/作业。后面章节基本都从这里展开。}
\exam{常考某功能属于哪个模块，如进程调度属于处理机管理，文件目录属于文件管理。}


\concept{操作系统提供的接口}
\begin{points}
  \item \term{命令接口}：用户直接输入命令控制系统，如 \code{ls}、\code{cd}、\code{mkdir}。Shell 本质是命令解释器。
  \item 命令接口分为联机用户接口和脱机用户接口。联机接口用于交互式控制；脱机接口常见于批处理作业控制语言。
  \item \term{程序接口}：程序通过系统调用请求 OS 服务，如 \code{open()}、\code{read()}、\code{fork()}。
  \item \term{图形接口}：窗口、菜单、鼠标等 GUI 形式，本质上也是调用系统服务。
  \item 广义指令常作为系统调用的另一种说法，属于程序接口。
\end{points}
\plain{人用命令/GUI 找 OS，程序用系统调用找 OS；危险操作必须让内核代办。}
\exam{常考命令接口、程序接口、GUI 的分类，以及系统调用是否属于程序接口。}

\concept{系统调用与特权操作}
\begin{points}
  \item I/O 指令、修改中断开关、修改页表、设置定时器等通常是特权指令，只能在内核态执行。
  \item 普通用户程序不能直接访问设备或随意操作内存，否则可能破坏其他进程或系统本身。
  \item 用户程序通过系统调用陷入内核，切换到内核态，由操作系统执行受保护的操作，再返回用户态。
\end{points}
\plain{人用命令/GUI 找 OS，程序用系统调用找 OS；危险操作必须让内核代办。}
\exam{常考命令接口、程序接口、GUI 的分类，以及系统调用是否属于程序接口。}


\section{操作系统结构}
\concept{常见结构}
\begin{points}
  \item \term{简单结构}：系统各部分边界不清，效率可能高，但维护困难。
  \item \term{层次结构}：从底层硬件到上层服务逐层构造，结构清晰，但层次划分和跨层调用可能带来开销。
  \item \term{微内核结构}：内核只保留最基本功能，如进程通信、低级地址空间管理、基本调度；文件系统、设备驱动等尽量放到用户态服务中。
  \item \term{模块化结构}：核心内核加可加载模块，兼顾性能与可扩展性。
  \item \term{虚拟机结构}：虚拟机监控器在一台物理机上抽象出多台虚拟计算机，每台虚拟机可运行自己的 OS。
\end{points}
\plain{OS 结构是在“好维护”和“高效率”之间权衡：简单结构快但乱，层次清楚但可能慢，微内核可靠但通信开销大。}
\exam{常考宏内核/微内核/模块化优缺点。}


\chapter{相关硬件基础}

\section{中央处理器}
\concept{CPU 的基本组成与工作流程}
\begin{points}
  \item CPU 的主要部件包括控制器、运算器、寄存器和高速缓存。
  \item CPU 的任务是执行指令，基本流程为：取指 $\rightarrow$ 译码 $\rightarrow$ 取操作数 $\rightarrow$ 执行指令。
  \item 影响 CPU 性能的因素包括主频、字长、指令集、高速缓存、核心数量、总线速度和处理技术。
\end{points}
\plain{硬件基础不用背太深，重点是知道哪些硬件机制支撑 OS：寄存器保存现场，PSW 保存状态，CPU 模式限制权限。}
\exam{常考 CPU 工作流程、单/多处理器、寄存器类别、PSW 中断位/状态位。}


\concept{单处理器与多处理器}
\begin{points}
  \item 单处理器系统只有一个运算处理器，同一时刻只能有一个进程在 CPU 上运行。
  \item 单 CPU 中，进程与进程只能并发不能并行；处理机与设备、设备与设备在硬件支持下可以并行工作。
  \item 多处理器系统有多个运算处理器，可分为对称多处理 SMP 和非对称多处理 AMP。
  \item SMP 中各处理器地位相同，共享内存和 I/O 系统；AMP 中主从处理器角色不同，主处理器负责调度与管理。
\end{points}
\plain{硬件基础不用背太深，重点是知道哪些硬件机制支撑 OS：寄存器保存现场，PSW 保存状态，CPU 模式限制权限。}
\exam{常考 CPU 工作流程、单/多处理器、寄存器类别、PSW 中断位/状态位。}


\concept{Flynn 分类}
\begin{points}
  \item SISD：单指令流单数据流，传统顺序单处理器。
  \item SIMD：单指令流多数据流，一条指令同时处理多个数据，如向量处理、GPU 中的某些模式。
  \item MISD：多指令流单数据流，实际很少见。
  \item MIMD：多指令流多数据流，多处理器/多核系统常见。
\end{points}
\plain{硬件基础不用背太深，重点是知道哪些硬件机制支撑 OS：寄存器保存现场，PSW 保存状态，CPU 模式限制权限。}
\exam{常考 CPU 工作流程、单/多处理器、寄存器类别、PSW 中断位/状态位。}


\section{寄存器、特权指令与处理器状态}
\concept{寄存器}
\begin{points}
  \item 寄存器是 CPU 内部高速存储单元，用于保存指令、地址、数据和状态信息。
  \item 常见寄存器包括程序计数器 PC、指令寄存器 IR、通用寄存器、栈指针 SP、程序状态字 PSW。
  \item 上下文切换时，操作系统需要保存当前进程的寄存器现场，并恢复下一个进程的现场。
\end{points}
\plain{硬件基础不用背太深，重点是知道哪些硬件机制支撑 OS：寄存器保存现场，PSW 保存状态，CPU 模式限制权限。}
\exam{常考 CPU 工作流程、单/多处理器、寄存器类别、PSW 中断位/状态位。}


\concept{特权指令与非特权指令}
\begin{points}
  \item 非特权指令可由用户程序直接执行，如普通算术、逻辑、数据移动指令。
  \item 特权指令只能由内核执行，如 I/O 指令、关中断、修改内存保护寄存器、修改页表、启动 DMA、设置定时器。
  \item 判断标准：凡是可能影响系统安全、影响其他进程、直接控制硬件资源的操作，通常属于特权操作。
\end{points}
\plain{普通程序不能直接碰全局危险资源；一旦需要 I/O、页表、中断等操作，就要进内核态。}
\exam{常考哪些指令是特权指令、系统调用是否引起模式切换、模式切换和进程切换区别。}


\concept{处理器状态：用户态与内核态}
\begin{points}
  \item 用户态下执行普通用户程序，只能执行非特权指令，不能直接访问受保护资源。
  \item 内核态下执行操作系统内核，可执行特权指令，访问核心数据结构。
  \item 从用户态到内核态的常见原因：系统调用、异常、外部中断。
  \item 从内核态返回用户态：中断/系统调用处理完成后，恢复用户进程现场并执行返回指令。
\end{points}
\plain{硬件基础不用背太深，重点是知道哪些硬件机制支撑 OS：寄存器保存现场，PSW 保存状态，CPU 模式限制权限。}
\exam{常考 CPU 工作流程、单/多处理器、寄存器类别、PSW 中断位/状态位。}


\concept{PSW 程序状态字}
\begin{points}
  \item PSW 保存 CPU 的运行状态信息，包括条件码、中断屏蔽位、处理器模式位等。
  \item 模式位用于区分用户态和内核态，是硬件支持操作系统保护机制的重要基础。
  \item 中断发生时，硬件通常会保存 PC/PSW 等关键信息，以便中断处理完成后回到原程序继续执行。
\end{points}
\plain{硬件基础不用背太深，重点是知道哪些硬件机制支撑 OS：寄存器保存现场，PSW 保存状态，CPU 模式限制权限。}
\exam{常考 CPU 工作流程、单/多处理器、寄存器类别、PSW 中断位/状态位。}


\section{中断处理}
\concept{中断的基本概念}
\begin{points}
  \item 中断是 CPU 暂停当前程序执行，转而执行相应处理程序，处理完成后再返回原程序的机制。
  \item 中断使 CPU 不必忙等 I/O 完成，是实现多道程序设计和 CPU/I/O 并行工作的关键。
  \item 中断来源包括外部设备中断、定时器中断、异常和系统调用陷入。
\end{points}
\plain{中断就是硬件或软件打断 CPU 当前执行，让 OS 先处理更紧急或必须处理的事件。}
\exam{常考中断处理步骤、中断与异常/系统调用区别、为什么中断能提高 CPU 与 I/O 并行度。}


\concept{中断处理过程}
\begin{points}
  \item 设备或内部事件提出中断请求。
  \item CPU 在适当时刻响应中断，保存断点和现场，切换到内核态。
  \item 根据中断向量找到对应的中断服务程序。
  \item 中断服务程序处理中断事件，例如读取设备状态、唤醒等待进程。
  \item 恢复被中断程序现场，返回原程序或转入调度程序。
\end{points}
\plain{中断就是硬件或软件打断 CPU 当前执行，让 OS 先处理更紧急或必须处理的事件。}
\exam{常考中断处理步骤、中断与异常/系统调用区别、为什么中断能提高 CPU 与 I/O 并行度。}


\concept{中断与异常、系统调用}
\begin{points}
  \item 外部中断通常来自 CPU 外部设备，如磁盘、网卡、键盘。
  \item 异常来自当前指令执行过程，如除零、非法指令、缺页异常。
  \item 系统调用是用户程序有意执行陷入指令，请求内核服务。
  \item 三者都会导致从用户态进入内核态，但触发原因不同。
\end{points}
\plain{人用命令/GUI 找 OS，程序用系统调用找 OS；危险操作必须让内核代办。}
\exam{常考命令接口、程序接口、GUI 的分类，以及系统调用是否属于程序接口。}


\section{操作系统的硬件支持}
\concept{内存保护}
\begin{points}
  \item 为防止进程访问非法内存，硬件需要支持地址变换和地址合法性检查。
  \item 简单系统中可用基址寄存器和界限寄存器检查地址是否越界。
  \item 分页/分段系统中由 MMU 根据页表或段表完成地址映射和保护检查。
\end{points}
\plain{把这个知识点放回本章主线理解：它回答的是 OS 为了管理资源、隐藏硬件、保证并发正确性而采用的某个对象、规则或步骤。}
\exam{常考选择/填空/简答，让你判断概念归属、比较相近概念，或按“定义-作用-机制-易错点”展开。}


\concept{定时器}
\begin{points}
  \item 定时器用于防止用户程序长期独占 CPU。
  \item 操作系统设置时间片，到时后定时器中断发生，CPU 转入内核执行调度程序。
  \item 分时系统和抢占式调度离不开定时器支持。
\end{points}
\plain{把这个知识点放回本章主线理解：它回答的是 OS 为了管理资源、隐藏硬件、保证并发正确性而采用的某个对象、规则或步骤。}
\exam{常考选择/填空/简答，让你判断概念归属、比较相近概念，或按“定义-作用-机制-易错点”展开。}


\concept{I/O 保护}
\begin{points}
  \item 用户程序不能直接执行 I/O 指令，必须通过系统调用请求内核完成。
  \item 这样可以统一设备分配、访问权限、缓冲、错误处理，防止多个进程直接抢设备。
  \item 因此 I/O 指令属于典型特权指令。
\end{points}
\plain{I/O 部分就是解决 CPU 快、设备慢、设备复杂且共享的问题。DMA/通道减少 CPU 搬运，缓冲平滑速度差，SPOOLing 把独占设备虚拟成共享设备。}
\exam{常考设备分类、四种 I/O 控制方式、DMA 与中断/通道区别、缓冲和 SPOOLing 的作用。}


\section{计算机启动与虚拟机}
\concept{计算机启动过程}
\begin{points}
  \item 上电后 CPU 从固定地址开始执行固件中的引导程序。
  \item 固件进行硬件自检和基本初始化，找到可启动设备。
  \item 引导程序把操作系统内核装入内存。
  \item 内核初始化数据结构、设备驱动、内存管理、进程管理，创建第一个用户态进程或系统服务。
\end{points}
\plain{开机不是直接进 OS，而是固件先跑，找到引导程序，再把 OS 内核装入内存。}
\exam{常考 BIOS/UEFI、引导程序、内核加载的大致顺序。}


\concept{虚拟机}
\begin{points}
  \item 虚拟机把底层物理机器抽象成多台逻辑机器，使多个 OS 能在同一物理平台上运行。
  \item 虚拟机监控器负责 CPU、内存、设备等资源的虚拟化与隔离。
  \item 直觉：普通 OS 把裸机虚拟成“适合应用使用的机器”；虚拟机监控器把裸机虚拟成“多台可运行 OS 的机器”。
\end{points}
\plain{不是资源真的变多，而是 OS 用映射、换入换出、排队等办法，让用户感觉自己有一份独立资源。}
\exam{常考虚拟处理机、虚拟存储器、虚拟设备的例子，以及虚拟性和资源复用的关系。}


\chapter{进程}

\section{进程引入}
\concept{前趋图}
\begin{points}
  \item 前趋图是有向无循环图，用于描述程序段、语句或进程之间的先后执行关系。
  \item 若 $P_i \rightarrow P_j$，表示 $P_i$ 必须在 $P_j$ 开始之前完成。
  \item 没有前趋的结点称为初始结点，没有后继的结点称为终止结点。
  \item 做题时根据箭头判断哪些操作必须先完成，哪些操作可以并发。
\end{points}
\plain{用有向无环图画出“谁必须在谁之前完成”，目的是识别哪些操作能并发。}
\exam{常考给前趋关系画图、数可并发的语句、写同步信号量约束。}


\concept{程序的顺序执行}
\begin{points}
  \item 顺序执行指一个程序中的多个操作严格按照规定顺序执行，例如输入 $I \rightarrow$ 计算 $C \rightarrow$ 输出 $P$。
  \item 特征一：顺序性，每一步必须在下一步开始前完成。
  \item 特征二：封闭性，程序运行结果不受外界因素影响。
  \item 特征三：可再现性，只要初始条件和执行环境相同，每次运行结果相同。
\end{points}
\plain{一个程序自己按固定步骤走，外界不插手，所以结果可重复。}
\exam{常考顺序性、封闭性、可再现性三个特征。}


\concept{程序的并发执行}
\begin{points}
  \item 并发执行指多个程序或程序段在时间上重叠推进，一个程序尚未结束，另一个程序已经开始。
  \item 特征一：间断性，程序表现为“执行 - 暂停 - 执行”。
  \item 特征二：失去封闭性，多个程序共享资源，资源状态可能被他人改变。
  \item 特征三：不可再现性，在初始条件相同的情况下，执行顺序不同也可能导致结果不同。
\end{points}
\plain{核心是“看起来同时推进”。单 CPU 上靠时间片交替，多 CPU 上才可能真的同一时刻运行。}
\exam{常考并发与并行区别、单处理机下最多几个进程运行、并发导致为什么需要同步互斥。}

\section{进程定义与描述}
\concept{进程的定义}
\begin{points}
  \item 进程是程序的一次执行过程，是系统进行资源分配和调度的基本单位之一。
  \item 进程不仅包含程序代码，还包含数据、堆栈、打开文件、寄存器现场、程序计数器和 PCB 等运行信息。
  \item 进程具有动态性、并发性、独立性、异步性和结构性。
\end{points}
\plain{程序是静态代码，进程是这段代码跑起来后的动态活动，还带资源和现场。}
\exam{常考程序与进程区别、进程实体由程序段/数据段/PCB 构成、进程是资源分配基本单位。}


\concept{程序与进程的区别}
\begin{points}
  \item 程序是静态的代码文件，进程是动态的执行活动。
  \item 一个程序可以对应多个进程，例如多个用户同时打开同一个编辑器。
  \item 一个进程在运行中可能执行多个程序片段，例如通过 \code{exec} 装入新程序。
  \item 程序没有生命周期状态，进程有创建、就绪、运行、阻塞、终止等状态。
\end{points}
\plain{程序是静态代码，进程是这段代码跑起来后的动态活动，还带资源和现场。}
\exam{常考程序与进程区别、进程实体由程序段/数据段/PCB 构成、进程是资源分配基本单位。}


\concept{PCB 进程控制块}
\begin{points}
  \item PCB 是操作系统感知和管理进程的唯一标志。
  \item PCB 保存进程标识信息、处理机状态、调度信息、进程控制信息、内存管理信息、I/O 状态信息等。
  \item 上下文切换本质上是保存当前进程 PCB 中的现场，并从下一个进程 PCB 中恢复现场。
  \item 选择题直觉：问“进程存在的唯一标志”通常答 PCB，不是程序、PID 或页表。
\end{points}
\plain{PCB 就是 OS 给每个进程建的档案，没有 PCB，OS 就没法调度和恢复这个进程。}
\exam{常考 PCB 内容、PCB 的作用、进程存在的唯一标志。}


\section{进程状态与转换}
\concept{三状态模型}
\begin{points}
  \item \term{就绪态}：进程已具备运行条件，只差 CPU。
  \item \term{运行态}：进程正在 CPU 上执行。
  \item \term{阻塞态/等待态}：进程因等待某事件不能运行，即使 CPU 空闲也不能执行。
  \item 状态转换：就绪 $\rightarrow$ 运行由调度引起；运行 $\rightarrow$ 就绪由时间片到或被抢占引起；运行 $\rightarrow$ 阻塞由等待 I/O/资源引起；阻塞 $\rightarrow$ 就绪由等待事件完成引起。
\end{points}
\plain{进程不是一直运行，它会在就绪、运行、阻塞等状态之间因调度或等待事件而移动。}
\exam{常考三/五/七状态转换图，尤其运行到阻塞、阻塞到就绪、就绪到运行的触发条件。}


\concept{五状态模型}
\begin{points}
  \item 在三状态基础上加入新建态和终止态。
  \item 新建态表示进程正在创建，尚未进入就绪队列。
  \item 终止态表示进程执行结束或被撤销，等待系统回收资源。
  \item 创建进程通常涉及分配 PCB、建立地址空间、装入程序、初始化资源、进入就绪队列。
\end{points}
\plain{进程不是一直运行，它会在就绪、运行、阻塞等状态之间因调度或等待事件而移动。}
\exam{常考三/五/七状态转换图，尤其运行到阻塞、阻塞到就绪、就绪到运行的触发条件。}


\concept{挂起状态}
\begin{points}
  \item 当内存紧张或系统需要调节多道程序度时，可把进程换出到外存，形成挂起状态。
  \item 就绪挂起：进程不在内存，但一旦调入内存即可运行。
  \item 阻塞挂起：进程不在内存，且还在等待某事件。
  \item 中级调度/交换调度常与挂起状态相关。
\end{points}
\plain{进程不是一直运行，它会在就绪、运行、阻塞等状态之间因调度或等待事件而移动。}
\exam{常考三/五/七状态转换图，尤其运行到阻塞、阻塞到就绪、就绪到运行的触发条件。}


\section{进程组织、控制与管理}
\concept{进程队列}
\begin{points}
  \item 就绪队列：所有驻留内存且处于就绪状态、等待 CPU 的进程集合。
  \item 设备队列：等待某个 I/O 设备的进程集合。
  \item 作业队列：外存上等待进入内存的作业集合。
  \item 进程在生命周期中会在多个队列之间迁移，调度程序就是从这些队列中选择合适对象。
\end{points}
\plain{队列是 OS 组织进程的方式：就绪的排就绪队列，等设备的排设备队列。}
\exam{常考进程在生命周期中如何在不同队列迁移。}


\concept{进程创建与终止}
\begin{points}
  \item 创建进程需要申请唯一标识、分配 PCB、分配必要资源、初始化上下文、插入就绪队列。
  \item 进程终止后，系统要回收内存、打开文件、I/O 资源、PCB 等。
  \item Unix/Linux 中常见系统调用：\code{fork()} 创建子进程，\code{exec()} 装入新程序，\code{wait()} 等待子进程，\code{exit()} 退出进程。
  \item \code{fork()} 返回两次：父进程中返回子进程 PID，子进程中返回 0；失败返回负值。
\end{points}
\plain{进程控制就是 OS 创建/撤销/阻塞/唤醒/切换进程，都会改 PCB 和队列。}
\exam{常考 fork/exec/wait/exit 以及父子进程数量。}


\concept{进程切换与模式切换}
\begin{points}
  \item 模式切换是用户态和内核态之间的切换，不一定更换正在运行的进程。
  \item 进程切换是 CPU 从一个进程切换到另一个进程，一定需要内核介入，通常伴随上下文保存和恢复。
  \item 系统调用会导致模式切换；是否导致进程切换，要看是否阻塞或调度。
  \item I/O 请求通常先从用户态陷入内核态，若需要等待 I/O 完成，则当前进程阻塞并发生进程切换。
\end{points}
\plain{模式切换只是换权限；进程切换还要保存/恢复另一个进程的现场，开销更大。}
\exam{常考系统调用是否必然进程切换、进程切换需要保存哪些上下文。}


\section{进程通信}
\concept{低级通信与高级通信}
\begin{points}
  \item 低级通信主要通过信号量等同步工具传递少量控制信息。
  \item 高级通信用于传递大量数据或结构化消息，常见方式包括共享存储、消息传递、管道。
\end{points}
\plain{进程地址空间相互隔离，所以要么共享一块区域，要么通过内核传消息。}
\exam{常考共享存储、消息传递、管道的特点和适用场景。}


\concept{共享存储}
\begin{points}
  \item 多个进程共享同一块内存区域，通过读写共享区交换数据。
  \item 优点：速度快，适合大量数据交换。
  \item 缺点：必须另行解决同步互斥问题，防止同时读写导致不一致。
\end{points}
\plain{多个进程都想用同一批资源，所以必须规定能不能一起用、谁先用、用完怎么释放。}
\exam{常考互斥共享/同时共享举例，或者问并发和共享为什么是 OS 最基本特征。}


\concept{消息传递}
\begin{points}
  \item 进程通过发送和接收消息通信，可由操作系统提供消息队列、邮箱等机制。
  \item 直接通信：发送者明确指定接收者。
  \item 间接通信：通过邮箱或消息队列进行。
  \item 优点：同步控制相对清晰，适合分布式系统；缺点是消息复制和内核介入可能带来开销。
\end{points}
\plain{进程地址空间相互隔离，所以要么共享一块区域，要么通过内核传消息。}
\exam{常考共享存储、消息传递、管道的特点和适用场景。}


\concept{管道}
\begin{points}
  \item 管道是一种半双工通信机制，常用于具有亲缘关系的进程之间。
  \item 匿名管道常用于父子进程；命名管道可用于无亲缘关系的进程。
  \item Shell 中的 \code{|} 就是把前一个命令的输出作为后一个命令的输入。
\end{points}
\plain{进程地址空间相互隔离，所以要么共享一块区域，要么通过内核传消息。}
\exam{常考共享存储、消息传递、管道的特点和适用场景。}


\chapter{线程}

\section{线程定义}
\concept{为什么引入线程}
\begin{points}
  \item 引入进程是为了描述和实现多道程序并发执行。
  \item 引入线程是为了减少并发执行的时空开销，让同一应用内部的多个任务能更轻量地并发。
  \item 早期进程同时承担“资源拥有单位”和“调度分派单位”两个角色，创建和切换开销较大。
  \item 线程把二者分开：进程更多作为资源拥有单位，线程作为 CPU 调度的基本执行单位。
\end{points}
\plain{线程是进程里面更轻的执行流，同进程线程共享资源，但各自有栈、寄存器和 PC。}
\exam{常考进程/线程区别、用户级线程和内核级线程优缺点、多对一/一对一/多对多模型。}


\concept{线程的定义}
\begin{points}
  \item 线程是进程内的一个执行单元，是进程内的可调度实体。
  \item 一个线程由线程 ID、程序计数器、寄存器集合和栈组成。
  \item 同一进程内的线程共享代码段、数据段、打开文件、地址空间等资源。
  \item 线程是“轻量级进程”，但不是独立资源拥有者。
\end{points}
\plain{线程是进程里面更轻的执行流，同进程线程共享资源，但各自有栈、寄存器和 PC。}
\exam{常考进程/线程区别、用户级线程和内核级线程优缺点、多对一/一对一/多对多模型。}


\concept{进程与线程区别}
\begin{points}
  \item 进程是资源分配的基本单位，线程是 CPU 调度的基本单位。
  \item 同一进程内线程共享地址空间，线程之间通信方便但相互影响更大。
  \item 不同进程地址空间隔离，通信开销较大但保护性更好。
  \item 线程创建、撤销、切换通常比进程轻量，因为不需要切换完整地址空间和大量资源。
\end{points}
\plain{程序是静态代码，进程是这段代码跑起来后的动态活动，还带资源和现场。}
\exam{常考程序与进程区别、进程实体由程序段/数据段/PCB 构成、进程是资源分配基本单位。}


\concept{多线程优点}
\begin{points}
  \item 响应性更好：例如界面线程响应用户，后台线程执行耗时任务。
  \item 资源共享方便：同一进程内线程天然共享内存和文件。
  \item 经济性更好：创建和切换开销较小。
  \item 可伸缩性更好：多核系统中多个线程可并行执行。
\end{points}
\plain{线程是进程里面更轻的执行流，同进程线程共享资源，但各自有栈、寄存器和 PC。}
\exam{常考进程/线程区别、用户级线程和内核级线程优缺点、多对一/一对一/多对多模型。}


\section{线程实现}
\concept{用户级线程}
\begin{points}
  \item 用户级线程由用户态线程库管理，内核不知道线程存在。
  \item 优点：创建和切换无需进入内核，速度快；调度策略可由应用定制。
  \item 缺点：一个线程执行阻塞系统调用，可能导致整个进程阻塞；无法真正利用多核并行。
\end{points}
\plain{线程是进程里面更轻的执行流，同进程线程共享资源，但各自有栈、寄存器和 PC。}
\exam{常考进程/线程区别、用户级线程和内核级线程优缺点、多对一/一对一/多对多模型。}


\concept{内核级线程}
\begin{points}
  \item 内核级线程由操作系统内核管理，调度器直接调度线程。
  \item 优点：一个线程阻塞不影响同进程其他线程；多线程可在多核上并行。
  \item 缺点：线程创建、切换、同步需要内核参与，开销比用户级线程大。
\end{points}
\plain{线程是进程里面更轻的执行流，同进程线程共享资源，但各自有栈、寄存器和 PC。}
\exam{常考进程/线程区别、用户级线程和内核级线程优缺点、多对一/一对一/多对多模型。}


\concept{多线程模型}
\begin{points}
  \item 多对一模型：多个用户线程映射到一个内核线程，管理轻量但阻塞问题明显。
  \item 一对一模型：每个用户线程映射到一个内核线程，并行能力强但内核线程数量受限。
  \item 多对多模型：多个用户线程映射到多个内核线程，兼顾灵活性与并行能力，但实现复杂。
\end{points}
\plain{线程是进程里面更轻的执行流，同进程线程共享资源，但各自有栈、寄存器和 PC。}
\exam{常考进程/线程区别、用户级线程和内核级线程优缺点、多对一/一对一/多对多模型。}


\concept{线程库}
\begin{points}
  \item 线程库向程序员提供创建、管理、同步线程的 API。
  \item 常见线程库包括 POSIX Pthreads、Windows threads、Java threads。
  \item 线程库可以完全在用户态实现，也可以作为内核级线程接口的包装。
\end{points}
\plain{线程是进程里面更轻的执行流，同进程线程共享资源，但各自有栈、寄存器和 PC。}
\exam{常考进程/线程区别、用户级线程和内核级线程优缺点、多对一/一对一/多对多模型。}


\concept{可重入代码}
\begin{points}
  \item 可重入代码指在被多个执行流同时调用或被中断后再次进入时，仍能得到正确结果的代码。
  \item 它通常不依赖不可保护的全局变量，不在函数内部保留共享静态状态，或对共享状态进行同步保护。
  \item 多线程、信号处理和中断处理程序常要求代码具备可重入性。
  \item 直觉：可重入代码像“每个人带自己的草稿纸”，不会抢同一张草稿纸写中间结果。
\end{points}
\plain{可重入代码像“公共只读函数”，多个线程同时进来也不会互相污染。}
\exam{常考可重入代码特征：不修改自身代码、不依赖全局可变状态或对其加保护。}


\chapter{进程调度}

\section{调度基本概念}
\concept{调度队列}
\begin{points}
  \item 作业队列：外存上等待进入系统的作业集合。
  \item 就绪队列：内存中已就绪、等待 CPU 的进程集合。
  \item 设备队列：等待某个 I/O 设备的进程集合。
  \item 进程在整个生命周期中会在这些队列之间迁移。
\end{points}
\plain{OS 把进程/作业放进不同队列，调度就是从队列里按规则挑下一个。}
\exam{常考作业队列/就绪队列/设备队列，高级/中级/低级调度分别管什么。}


\concept{三级调度}
\begin{points}
  \item 高级调度又称作业调度/长程调度，决定哪些作业从外存进入内存，影响多道程序度。
  \item 中级调度又称交换调度，决定哪些进程换入、换出内存，用于缓解内存紧张。
  \item 低级调度又称进程调度/短程调度，决定就绪队列中哪个进程获得 CPU，运行频率最高。
  \item 记忆：高级管“进不进系统”，中级管“在不在内存”，低级管“谁上 CPU”。
\end{points}
\plain{OS 把进程/作业放进不同队列，调度就是从队列里按规则挑下一个。}
\exam{常考作业队列/就绪队列/设备队列，高级/中级/低级调度分别管什么。}


\concept{CPU 密集型与 I/O 密集型}
\begin{points}
  \item CPU 密集型作业计算时间多，CPU burst 长，I/O 次数相对少。
  \item I/O 密集型作业 I/O 操作多，CPU burst 短。
  \item 作业调度应合理搭配两类作业：CPU 密集型过多会导致 I/O 设备空闲；I/O 密集型过多会导致就绪队列空、CPU 空闲。
\end{points}
\plain{硬件基础不用背太深，重点是知道哪些硬件机制支撑 OS：寄存器保存现场，PSW 保存状态，CPU 模式限制权限。}
\exam{常考 CPU 工作流程、单/多处理器、寄存器类别、PSW 中断位/状态位。}


\concept{调度程序与分派程序}
\begin{points}
  \item 调度程序负责从就绪队列中选择下一个运行进程。
  \item 分派程序负责完成上下文切换、切换到用户态、跳转到被选进程的正确位置。
  \item 分派延迟是停止一个进程并启动另一个进程所需时间，应尽可能短。
\end{points}
\plain{把这个知识点放回本章主线理解：它回答的是 OS 为了管理资源、隐藏硬件、保证并发正确性而采用的某个对象、规则或步骤。}
\exam{常考选择/填空/简答，让你判断概念归属、比较相近概念，或按“定义-作用-机制-易错点”展开。}


\section{调度目标与评价指标}
\concept{常见评价指标}
\begin{points}
  \item CPU 利用率：CPU 忙碌时间比例，越高越好。
  \item 吞吐量：单位时间完成的进程数量，越高越好。
  \item 周转时间：进程从提交到完成的总时间，越短越好。
  \item 等待时间：进程在就绪队列中等待 CPU 的总时间，越短越好。
  \item 响应时间：从提交请求到首次响应的时间，分时系统尤其关注。
\end{points}
\plain{调度题本质是按时间轴算完成时间，再由完成时间反推周转、等待、响应。}
\exam{常考平均周转时间、带权周转时间、等待时间计算，务必画甘特图。}


\concept{周转时间公式}
\begin{points}
  \item 完成时间 $C_i$，到达时间 $A_i$，服务时间 $S_i$。
  \item 周转时间 $T_i=C_i-A_i$。
  \item 带权周转时间 $W_i=\dfrac{T_i}{S_i}$。
  \item 平均周转时间和平均带权周转时间就是对所有进程取平均。
\end{points}
\plain{调度题本质是按时间轴算完成时间，再由完成时间反推周转、等待、响应。}
\exam{常考平均周转时间、带权周转时间、等待时间计算，务必画甘特图。}


\section{调度方式}
\concept{抢占式与非抢占式}
\begin{points}
  \item 非抢占式调度：进程一旦获得 CPU，除非主动放弃、阻塞或结束，否则不会被强行剥夺。
  \item 抢占式调度：操作系统可因时间片到、更高优先级进程到达等原因剥夺 CPU。
  \item 分时系统通常需要抢占式调度；批处理系统可以使用非抢占式调度。
\end{points}
\plain{抢占就是正在运行的进程也可能被更高优先级/时间片到期打断。}
\exam{常考抢占式与非抢占式适用场景。}


\section{常用调度算法}
\concept{先来先服务 FCFS}
\begin{points}
  \item 按进程到达就绪队列的先后顺序分配 CPU。
  \item 优点：简单、公平、不会饥饿。
  \item 缺点：短作业可能排在长作业后面，平均等待时间可能很长，存在“护航效应”。
  \item 适合批处理场景，不适合交互式短任务较多的场景。
\end{points}
\plain{谁先来谁先上，最朴素但容易让短作业排在长作业后面。}
\exam{常考护航效应和平均等待时间计算。}


\concept{短作业优先 SJF 与最短剩余时间优先 SRTF}
\begin{points}
  \item SJF 选择预计运行时间最短的作业先执行，是非抢占式算法。
  \item SRTF 是抢占式版本，总是选择剩余运行时间最短的进程。
  \item 优点：理论上可使平均等待时间最短。
  \item 缺点：需要预测运行时间；长作业可能长期等待，产生饥饿。
\end{points}
\plain{优先做短的，平均等待时间好看，但长作业可能一直等。}
\exam{常考非抢占 SJF 和抢占式 SRTF 的甘特图区别。}


\concept{优先级调度}
\begin{points}
  \item 为每个进程赋予优先级，优先级高者先运行。
  \item 可分为抢占式和非抢占式。
  \item 静态优先级创建后不变，动态优先级可根据等待时间、资源使用情况改变。
  \item 缺点是低优先级进程可能饥饿；常用老化 aging 提高等待过久进程的优先级。
\end{points}
\plain{给进程分等级，重要的先跑；问题是低优先级可能饥饿。}
\exam{常考静态/动态优先级、抢占/非抢占、老化技术防饥饿。}


\concept{时间片轮转 RR}
\begin{points}
  \item 每个就绪进程轮流获得一个时间片，时间片用完仍未完成则回到就绪队列尾部。
  \item 适合分时系统，强调响应时间和公平性。
  \item 时间片太大，退化为 FCFS；时间片太小，上下文切换开销过大。
  \item 做题时严格画时间轴，注意新到达进程插入队列的位置和时间片剩余规则。
\end{points}
\plain{每人轮流跑一小段，适合分时系统；时间片太大像 FCFS，太小切换开销大。}
\exam{常考时间片轮转甘特图，以及时间片大小对响应时间和开销的影响。}


\concept{高响应比优先 HRRN}
\begin{points}
  \item 响应比公式：$R=\dfrac{等待时间+服务时间}{服务时间}=1+\dfrac{等待时间}{服务时间}$。
  \item 等待时间越长，响应比越大；服务时间越短，响应比也越大。
  \item 该算法兼顾短作业优先和防止长作业饥饿。
  \item 通常是非抢占式调度，用在作业调度中较常见。
\end{points}
\plain{每人轮流跑一小段，适合分时系统；时间片太大像 FCFS，太小切换开销大。}
\exam{常考时间片轮转甘特图，以及时间片大小对响应时间和开销的影响。}


\concept{多级队列调度}
\begin{points}
  \item 将就绪队列划分为多个队列，不同队列服务不同类型进程，如系统进程、交互进程、批处理进程。
  \item 队列之间有固定优先级或时间比例分配，队列内部可使用不同调度算法。
  \item 固定多级队列的问题是低优先级队列可能长期得不到服务。
\end{points}
\plain{新来的先给高优先级小时间片；用得越久越往下掉，适合交互和批处理混合。}
\exam{常考多级反馈队列调度过程，注意抢占、降级、时间片没用完是否保留。}


\concept{多级反馈队列 MLFQ}
\begin{points}
  \item 多级反馈队列允许进程在不同优先级队列之间移动。
  \item 新进程通常先进入高优先级队列，时间片较短，若用完时间片还未完成则降级。
  \item 等待过久的低优先级进程可被提升，避免饥饿。
  \item 直觉：先假设新来的任务可能是短交互任务，给它快速响应；若它表现得像长计算任务，就逐步降到低优先级。
\end{points}
\plain{队列是 OS 组织进程的方式：就绪的排就绪队列，等设备的排设备队列。}
\exam{常考进程在生命周期中如何在不同队列迁移。}


\section{实时调度、多处理器调度与线程调度}
\concept{实时调度基本概念}
\begin{points}
  \item 实时系统中，正确结果如果迟到，也可能等同于错误结果。
  \item 周期性任务常用处理时间 $t_i$、截止期限 $d_i$、周期 $p_i$ 描述，通常 $0\le t_i\le d_i\le p_i$。
  \item 必要的可调度条件之一：$\sum_{i=1}^{m}\dfrac{t_i}{p_i}\le 1$，表示 CPU 总利用需求不能超过 100\%。
\end{points}
\plain{实时调度不只看快慢，而是看能否在截止期限前完成。}
\exam{常考硬/软实时、中断延迟/分派延迟、RMS/EDF 和可调度条件。}


\concept{速率单调 RMS 与最早截止期限 EDF}
\begin{points}
  \item RMS 是静态优先级算法，周期越短，优先级越高。
  \item EDF 是动态优先级算法，截止期限越近，优先级越高。
  \item EDF 在单处理器、理想条件下可达到更高利用率，但实现和开销更复杂。
\end{points}
\plain{实时调度不只看快慢，而是看能否在截止期限前完成。}
\exam{常考硬/软实时、中断延迟/分派延迟、RMS/EDF 和可调度条件。}


\concept{多处理器调度}
\begin{points}
  \item 非对称多处理调度中，一个主处理器负责调度，其他处理器执行任务。
  \item 对称多处理调度中，各处理器都可参与调度，常见挑战是负载均衡和缓存亲和性。
  \item 处理器亲和性指尽量让进程继续在原 CPU 上运行，以利用缓存中的数据。
\end{points}
\plain{硬件基础不用背太深，重点是知道哪些硬件机制支撑 OS：寄存器保存现场，PSW 保存状态，CPU 模式限制权限。}
\exam{常考 CPU 工作流程、单/多处理器、寄存器类别、PSW 中断位/状态位。}


\concept{线程调度}
\begin{points}
  \item 用户级线程由线程库调度，内核只看到进程。
  \item 内核级线程由内核调度，可以在多个处理器上并行。
  \item 线程调度更强调轻量切换和同一进程内线程之间的同步关系。
\end{points}
\plain{线程是进程里面更轻的执行流，同进程线程共享资源，但各自有栈、寄存器和 PC。}
\exam{常考进程/线程区别、用户级线程和内核级线程优缺点、多对一/一对一/多对多模型。}


\chapter{进程同步}

\section{同步与互斥的概念}
\concept{并发进程的制约关系}
\begin{points}
  \item 多道程序系统中，由于资源共享或进程合作，并发进程之间形成相互制约关系。
  \item 间接相互制约：多个进程竞争同一资源，例如多个进程争用打印机。
  \item 直接相互制约：进程之间存在合作顺序，例如生产者生产后消费者才能消费。
  \item 进程同步的任务是协调这些关系，使并发进程能正确共享资源和相互合作，保证结果可再现。
\end{points}
\plain{核心是“看起来同时推进”。单 CPU 上靠时间片交替，多 CPU 上才可能真的同一时刻运行。}
\exam{常考并发与并行区别、单处理机下最多几个进程运行、并发导致为什么需要同步互斥。}


\concept{与时间有关的错误/race condition}
\begin{points}
  \item 多个进程并发读写共享变量，执行次序不同可能导致不同结果，这种错误称为与时间有关的错误或竞争条件。
  \item 例子：订票系统中 A、B 同时读取剩余票数 $x=1$，都判断有票并出票，最终可能超卖。
  \item 根本原因：对共享变量的复合操作没有互斥保护。
  \item 解决思路：把读-改-写共享变量的代码段放入临界区，一次只允许一个进程进入。
\end{points}
\plain{并发程序共享数据时，执行顺序一变结果就变，所以要把关键代码段保护起来。}
\exam{常考直接/间接制约、竞争条件、临界资源/临界区、四条临界区准则。}


\concept{临界资源与临界区}
\begin{points}
  \item 临界资源是一段时间内只允许一个进程使用的资源。
  \item 临界区是进程中访问临界资源的代码段。
  \item 进入区负责申请进入临界区；退出区负责释放临界区；剩余区是其他代码。
  \item 互斥要求：任意时刻最多一个进程在相关临界区内执行。
\end{points}
\plain{并发程序共享数据时，执行顺序一变结果就变，所以要把关键代码段保护起来。}
\exam{常考直接/间接制约、竞争条件、临界资源/临界区、四条临界区准则。}


\concept{临界区访问原则}
\begin{points}
  \item \key{空闲让进}：临界区空闲时，应允许一个请求进入的进程立即进入。
  \item \key{忙则等待}：临界区已有进程时，其他请求进程必须等待。
  \item \key{有限等待}：等待进入临界区的进程不能无限期等待，避免饥饿。
  \item \key{让权等待}：不能进入临界区的进程应释放处理机，不应长时间忙等浪费 CPU。
\end{points}
\plain{并发程序共享数据时，执行顺序一变结果就变，所以要把关键代码段保护起来。}
\exam{常考直接/间接制约、竞争条件、临界资源/临界区、四条临界区准则。}

\section{互斥的实现方法}
\concept{软件方法}
\begin{points}
  \item 单标志法、双标志先检查法、双标志后检查法、Peterson 算法都是早期软件互斥方法。
  \item 软件方法的关键是用共享变量表达“谁想进”和“轮到谁进”。
  \item Peterson 算法能满足互斥、空闲让进和有限等待，但主要适合两个进程的理论模型。
  \item 软件方法缺点是复杂、容易出错、依赖忙等，难以推广。
\end{points}
\plain{互斥可以靠软件协议，也可以靠关中断/测试并设置等硬件原子指令；现代更常封装成锁或信号量。}
\exam{常考忙等缺点、关中断是否适合多处理器、Test-and-Set 的作用。}


\concept{硬件方法}
\begin{points}
  \item 关中断：单处理器中进入临界区前关中断，退出后开中断，可防止进程被中断切换。
  \item Test-and-Set 指令：原子地测试并设置锁变量。
  \item Swap 指令：原子地交换两个变量的值。
  \item 硬件方法可保证互斥，但若采用自旋锁，会造成忙等，不满足让权等待。
\end{points}
\plain{互斥可以靠软件协议，也可以靠关中断/测试并设置等硬件原子指令；现代更常封装成锁或信号量。}
\exam{常考忙等缺点、关中断是否适合多处理器、Test-and-Set 的作用。}


\section{信号量与 P/V 操作}
\concept{信号量定义}
\begin{points}
  \item 信号量 semaphore 由整数值 $s$ 和等待队列 $q$ 组成。
  \item 除初始化外，信号量的值只能通过 P 操作和 V 操作改变。
  \item 信号量 $s>0$ 时，表示可用资源数；$s=0$ 表示无可用资源但无人等待；$s<0$ 时，$|s|$ 表示等待队列中的进程数。
  \item P/V 操作必须是原子操作，执行过程中不能被其他进程打断。
\end{points}
\plain{P 是先占/先减，不够就睡；V 是释放/先加，如果有人等就叫醒一个。}
\exam{常考 P/V 原语含义、s>0/s<0 的解释、互斥信号量和资源信号量的初值。}


\concept{P 操作}
\begin{points}
  \item P(S)：先执行 $s=s-1$。
  \item 若减完后 $s<0$，说明资源不够，当前进程阻塞，进入该信号量等待队列，并转调度程序。
  \item 若减完后 $s\ge 0$，说明成功获得资源，进程继续执行。
  \item 直觉：P 是“申请/占用一个资源”，先登记需求，再判断是否需要等待。
\end{points}
\plain{P 是先占/先减，不够就睡；V 是释放/先加，如果有人等就叫醒一个。}
\exam{常考 P/V 原语含义、s>0/s<0 的解释、互斥信号量和资源信号量的初值。}


\concept{V 操作}
\begin{points}
  \item V(S)：先执行 $s=s+1$。
  \item 若加完后 $s\le 0$，说明仍有进程在等待队列中，唤醒一个等待进程并插入就绪队列。
  \item 若加完后 $s>0$，说明没有进程等待，只是增加了一个可用资源。
  \item 直觉：V 是“释放一个资源”。即使加完后 $s\le 0$，也要唤醒，因为负数表示有等待者；释放的资源应交给其中一个等待进程。
\end{points}
\plain{P 是先占/先减，不够就睡；V 是释放/先加，如果有人等就叫醒一个。}
\exam{常考 P/V 原语含义、s>0/s<0 的解释、互斥信号量和资源信号量的初值。}

\concept{整型信号量与记录型信号量}
\begin{points}
  \item 整型信号量只用整数表示资源数，P 操作中可能忙等。
  \item 记录型信号量增加等待队列，资源不足时让进程阻塞，满足让权等待。
  \item 考试中常默认 P/V 是记录型信号量，阻塞和唤醒由 OS 完成。
\end{points}
\plain{P 是先占/先减，不够就睡；V 是释放/先加，如果有人等就叫醒一个。}
\exam{常考 P/V 原语含义、s>0/s<0 的解释、互斥信号量和资源信号量的初值。}


\section{经典同步问题}
\concept{生产者-消费者问题}
\begin{points}
  \item 问题：多个生产者向有界缓冲区放产品，多个消费者从缓冲区取产品。
  \item 需要三个信号量：\code{mutex=1} 保护缓冲区互斥访问，\code{empty=n} 表示空缓冲区数，\code{full=0} 表示满缓冲区数。
  \item 生产者顺序：生产产品 $\rightarrow$ P(empty) $\rightarrow$ P(mutex) $\rightarrow$ 放入缓冲区 $\rightarrow$ V(mutex) $\rightarrow$ V(full)。
  \item 消费者顺序：P(full) $\rightarrow$ P(mutex) $\rightarrow$ 取出产品 $\rightarrow$ V(mutex) $\rightarrow$ V(empty) $\rightarrow$ 消费产品。
  \item 注意：先申请同步资源 empty/full，再申请互斥 mutex，避免拿着锁等待条件导致其他进程无法改变条件。
\end{points}
\plain{经典同步题先找资源数量，再找互斥，再找先后顺序。}
\exam{常考补全 P/V 代码、判断会不会死锁、修改信号量初值。}


\concept{读者-写者问题}
\begin{points}
  \item 多个读者可同时读共享文件，但写者必须互斥写；写者与读者也互斥。
  \item 读者优先方案中，第一个读者负责阻止写者，最后一个读者负责释放写者。
  \item 需要计数变量 \code{readcount} 记录当前读者数，并用互斥信号量保护它。
  \item 易错点：读者之间不是互斥读文件，而是互斥修改 \code{readcount}。
\end{points}
\plain{经典同步题先找资源数量，再找互斥，再找先后顺序。}
\exam{常考补全 P/V 代码、判断会不会死锁、修改信号量初值。}


\concept{哲学家进餐问题}
\begin{points}
  \item 五个哲学家围桌，每人左右各一只筷子，必须同时拿到两只筷子才能进餐。
  \item 若所有哲学家都先拿左筷子再等右筷子，可能形成死锁。
  \item 解决方法包括：最多允许四人同时取筷子；奇偶哲学家取筷子顺序相反；一次性申请两只筷子；使用管程。
  \item 核心直觉：避免每个进程“占有一个资源，同时等待另一个资源”。
\end{points}
\plain{经典同步题先找资源数量，再找互斥，再找先后顺序。}
\exam{常考补全 P/V 代码、判断会不会死锁、修改信号量初值。}


\section{信号量集机制}
\concept{AND 型信号量}
\begin{points}
  \item AND 型信号量适用于进程同时需要多类资源、每类一个的情况。
  \item 基本思想：进程需要的多类资源一次性全部分配；只要有一个资源不满足，其他资源也不分配。
  \item 这样可以避免进程占有部分资源再等待其他资源，从而降低死锁风险。
  \item P 原语称为 SP/Swait，V 原语称为 SV/Ssignal。
\end{points}
\plain{P 是先占/先减，不够就睡；V 是释放/先加，如果有人等就叫醒一个。}
\exam{常考 P/V 原语含义、s>0/s<0 的解释、互斥信号量和资源信号量的初值。}


\concept{一般信号量集}
\begin{points}
  \item 一般信号量集是 AND 型的扩展：一次可申请多类资源，每类资源可申请多个。
  \item 当某类资源低于下限或不能满足申请量时，本次原语不分配任何资源。
  \item 它把“多个 P 操作”合并成一个原子判断，避免中间状态。
\end{points}
\plain{P 是先占/先减，不够就睡；V 是释放/先加，如果有人等就叫醒一个。}
\exam{常考 P/V 原语含义、s>0/s<0 的解释、互斥信号量和资源信号量的初值。}


\section{管程}
\concept{管程是什么}
\begin{points}
  \item 管程是一种高级同步机制，把共享数据、对共享数据的操作、同步条件封装在一个模块中。
  \item 管程保证任意时刻只有一个进程/线程在管程内部执行，因此天然提供互斥。
  \item 条件变量用于表达“某个条件不满足时等待，条件满足时被唤醒”。
  \item 典型操作包括 \code{wait} 和 \code{signal}。
\end{points}
\plain{管程把共享数据和操作封装起来，自动保证一次只有一个进程进管程。}
\exam{常考管程和 P/V 区别：管程是高级语言级封装，P/V 更底层更灵活也更易错。}


\concept{管程与 P/V 的区别}
\begin{points}
  \item P/V 是低级同步原语，程序员需要自己安排每个 P、V 的位置，容易写错。
  \item 管程把共享变量和互斥规则封装起来，进入管程自动互斥，程序员主要表达条件等待。
  \item P/V 更像“手动锁门开门”；管程更像“把房间设计成一次只能进一个人，并提供排队等待条件”。
  \item 管程更结构化，适合高级语言；信号量更基础，适合描述底层同步。
\end{points}
\plain{管程把共享数据和操作封装起来，自动保证一次只有一个进程进管程。}
\exam{常考管程和 P/V 区别：管程是高级语言级封装，P/V 更底层更灵活也更易错。}


\chapter{死锁}

\section{死锁概念}
\concept{死锁定义}
\begin{points}
  \item 死锁是多个进程因竞争系统资源而造成的一种僵局，若无外力作用，这些进程都将永远不能向前推进。
  \item 若进程集合中的每个进程都在等待只能由该集合中其他进程引发的事件，则该集合发生死锁。
  \item 典型例子：P1 占有打印机等待读卡机，P2 占有读卡机等待打印机，双方都无法继续。
\end{points}
\plain{死锁就是大家都拿着别人需要的资源不放，结果谁都走不下去。}
\exam{常考四个必要条件、预防/避免/检测解除区别、银行家算法安全序列。}


\concept{死锁与饥饿、死循环}
\begin{points}
  \item 死锁：一组进程相互等待，若无外力永远不能推进。
  \item 饥饿：某进程长期得不到所需资源，但系统整体仍在推进。
  \item 死循环：一个进程自身逻辑错误导致反复执行，不一定涉及资源等待。
  \item 活锁：进程没有阻塞，但不断改变状态、相互让步，导致没有实际进展。
\end{points}
\plain{死锁就是大家都拿着别人需要的资源不放，结果谁都走不下去。}
\exam{常考四个必要条件、预防/避免/检测解除区别、银行家算法安全序列。}


\section{资源分类与死锁必要条件}
\concept{资源分类}
\begin{points}
  \item 可剥夺资源：可从进程手中强行取回且不破坏计算，例如 CPU、内存页框。
  \item 不可剥夺资源：一旦分配，不能强行剥夺，否则会造成错误，例如打印机、部分外设。
  \item 可重用资源：使用完可再次分配，如设备、内存。
  \item 消耗性资源：使用后消失，如消息、信号。
  \item 死锁主要与不可剥夺的可重用资源有关。
\end{points}
\plain{把这个知识点放回本章主线理解：它回答的是 OS 为了管理资源、隐藏硬件、保证并发正确性而采用的某个对象、规则或步骤。}
\exam{常考选择/填空/简答，让你判断概念归属、比较相近概念，或按“定义-作用-机制-易错点”展开。}


\concept{死锁产生的四个必要条件}
\begin{points}
  \item \key{互斥条件}：资源一次只能被一个进程使用。
  \item \key{请求和保持条件}：进程已保持至少一个资源，同时又请求新的资源。
  \item \key{不可剥夺条件}：资源不能被强行夺走，只能由占有进程主动释放。
  \item \key{循环等待条件}：存在进程资源等待环，$P_1$ 等 $P_2$，$P_2$ 等 $P_3$，$\ldots$，$P_n$ 等 $P_1$。
\end{points}
\plain{死锁就是大家都拿着别人需要的资源不放，结果谁都走不下去。}
\exam{常考四个必要条件、预防/避免/检测解除区别、银行家算法安全序列。}

\section{处理死锁的基本方法}
\concept{四类方法总览}
\begin{points}
  \item 死锁预防：通过制度性限制破坏死锁必要条件之一。
  \item 死锁避免：运行时根据资源需求判断分配是否安全，典型算法是银行家算法。
  \item 死锁检测与解除：允许死锁发生，定期检测，发生后恢复。
  \item 鸵鸟策略：忽略死锁问题，适用于死锁概率很低且处理成本过高的系统。
\end{points}
\plain{把这个知识点放回本章主线理解：它回答的是 OS 为了管理资源、隐藏硬件、保证并发正确性而采用的某个对象、规则或步骤。}
\exam{常考选择/填空/简答，让你判断概念归属、比较相近概念，或按“定义-作用-机制-易错点”展开。}


\concept{死锁预防}
\begin{points}
  \item 破坏互斥条件：把独占设备改造成共享设备，例如 SPOOLing 把打印机虚拟成共享设备。但并非所有资源都能共享。
  \item 破坏请求和保持条件：进程运行前一次性申请所有资源，或申请新资源前释放已占资源。缺点是资源利用率低、可能饥饿。
  \item 破坏不可剥夺条件：资源申请失败时剥夺其已占资源。适用于 CPU、内存等可剥夺资源，不适合打印机等。
  \item 破坏循环等待条件：对资源统一编号，进程按编号递增顺序申请资源。
\end{points}
\plain{死锁就是大家都拿着别人需要的资源不放，结果谁都走不下去。}
\exam{常考四个必要条件、预防/避免/检测解除区别、银行家算法安全序列。}


\concept{安全状态与不安全状态}
\begin{points}
  \item 安全状态指系统能找到一个安全序列，使所有进程都能按某种顺序获得所需资源并完成。
  \item 不安全状态不一定已经死锁，但可能发展为死锁。
  \item 死锁状态一定是不安全状态；不安全状态不一定是死锁状态。
\end{points}
\plain{进程不是一直运行，它会在就绪、运行、阻塞等状态之间因调度或等待事件而移动。}
\exam{常考三/五/七状态转换图，尤其运行到阻塞、阻塞到就绪、就绪到运行的触发条件。}


\concept{银行家算法}
\begin{points}
  \item 银行家算法用于死锁避免，假设系统知道每个进程的最大资源需求。
  \item 关键数据结构：Available 可用资源，Max 最大需求，Allocation 已分配，Need 尚需，且 $Need=Max-Allocation$。
  \item 资源请求满足 $Request_i\le Need_i$ 且 $Request_i\le Available$ 后，先试探性分配。
  \item 再执行安全性检查：寻找一个进程满足 $Need_i\le Work$，假设其完成并释放资源，更新 $Work=Work+Allocation_i$，直到所有进程完成或找不到下一个。
  \item 若试分配后仍安全，则真正分配；否则让进程等待。
\end{points}
\plain{死锁就是大家都拿着别人需要的资源不放，结果谁都走不下去。}
\exam{常考四个必要条件、预防/避免/检测解除区别、银行家算法安全序列。}


\concept{资源分配图}
\begin{points}
  \item 资源分配图用圆表示进程，用方框表示资源类型，方框内点表示资源实例。
  \item 请求边：进程 $\rightarrow$ 资源；分配边：资源实例 $\rightarrow$ 进程。
  \item 若每类资源只有一个实例，图中有环等价于发生死锁。
  \item 若资源类型有多个实例，有环只是死锁的必要条件，不一定死锁。
\end{points}
\plain{死锁就是大家都拿着别人需要的资源不放，结果谁都走不下去。}
\exam{常考四个必要条件、预防/避免/检测解除区别、银行家算法安全序列。}


\concept{死锁检测与解除}
\begin{points}
  \item 死锁检测允许资源正常分配，周期性检查是否存在死锁。
  \item 检测后可通过撤销进程、回滚进程、剥夺资源等方式解除死锁。
  \item 选择牺牲进程时可考虑优先级、已运行时间、已占资源、还需资源、回滚代价等。
  \item 缺点是检测和恢复都有代价，且可能造成进程饥饿。
\end{points}
\plain{死锁就是大家都拿着别人需要的资源不放，结果谁都走不下去。}
\exam{常考四个必要条件、预防/避免/检测解除区别、银行家算法安全序列。}


\chapter{存储器管理}

\section{存储器管理基本概念}
\concept{管理对象与目标}
\begin{points}
  \item 存储管理的对象是主存储器，即内存。
  \item 主存通常分为系统区和用户区：系统区存放内核程序与数据结构，用户区存放用户进程代码和数据。
  \item 目标：为用户提供方便、安全、充分大的存储空间，支持大型应用和系统程序运行。
\end{points}
\plain{把这个知识点放回本章主线理解：它回答的是 OS 为了管理资源、隐藏硬件、保证并发正确性而采用的某个对象、规则或步骤。}
\exam{常考选择/填空/简答，让你判断概念归属、比较相近概念，或按“定义-作用-机制-易错点”展开。}


\concept{存储器管理四个功能}
\begin{points}
  \item 分配与回收：为进程分配内存，进程结束或换出时回收。
  \item 抽象与映射：让进程看到连续的逻辑地址空间，并把逻辑地址转换成物理地址。
  \item 保护与共享：防止进程非法访问他人空间，同时允许受控共享。
  \item 存储扩充：在逻辑上提供比实际物理内存更大的存储空间。
\end{points}
\plain{把这个知识点放回本章主线理解：它回答的是 OS 为了管理资源、隐藏硬件、保证并发正确性而采用的某个对象、规则或步骤。}
\exam{常考选择/填空/简答，让你判断概念归属、比较相近概念，或按“定义-作用-机制-易错点”展开。}


\concept{逻辑地址与物理地址}
\begin{points}
  \item 逻辑地址是 CPU 运行的进程所看到的地址，也叫虚地址或相对地址。
  \item 物理地址是内存中的实际地址，也叫实地址或绝对地址。
  \item 地址变换/地址映射/重定位是把逻辑地址转换为物理地址的过程。
  \item MMU 是常见的硬件支持，负责地址转换和合法性检查。
\end{points}
\plain{程序看到的是逻辑地址，内存条上真正用的是物理地址，中间靠装载器/MMU 转换。}
\exam{常考逻辑地址到物理地址转换、静态/动态重定位、越界检查。}


\concept{地址绑定时机}
\begin{points}
  \item 编译时绑定：若装入地址已知，编译器可直接生成绝对地址；位置改变则需重新编译。
  \item 装入时绑定：编译生成可重定位代码，装入时确定物理地址。
  \item 运行时绑定：进程运行中仍可移动，地址在执行时由硬件转换，现代分页/分段常用。
\end{points}
\plain{程序看到的是逻辑地址，内存条上真正用的是物理地址，中间靠装载器/MMU 转换。}
\exam{常考逻辑地址到物理地址转换、静态/动态重定位、越界检查。}


\concept{静态重定位与动态重定位}
\begin{points}
  \item 静态重定位在程序装入内存时一次性完成地址修改，运行中不能随意移动。
  \item 动态重定位在程序运行时进行，每次访问内存时由重定位寄存器或 MMU 转换。
  \item 动态重定位支持进程换入换出、移动和虚拟存储，但需要硬件支持。
\end{points}
\plain{程序看到的是逻辑地址，内存条上真正用的是物理地址，中间靠装载器/MMU 转换。}
\exam{常考逻辑地址到物理地址转换、静态/动态重定位、越界检查。}


\section{连续分配管理}
\concept{单一连续分配}
\begin{points}
  \item 内存用户区一次只装入一道用户程序，用户程序独占用户区。
  \item 实现简单，适合早期单道系统。
  \item 缺点是内存利用率低，不能支持多道程序并发。
\end{points}
\plain{连续分配要求一整块内存，容易产生碎片；分页把连续要求打碎。}
\exam{常考内部/外部碎片区别、首次/最佳/最坏适应算法、伙伴系统合并。}


\concept{固定分区分配}
\begin{points}
  \item 系统启动时把内存划分为若干固定大小分区，每个分区装入一个进程。
  \item 分区大小可相等，也可不等。
  \item 优点：实现简单，支持多道程序。
  \item 缺点：分区内部未用空间形成内部碎片；分区数限制多道程序道数。
\end{points}
\plain{连续分配要求一整块内存，容易产生碎片；分页把连续要求打碎。}
\exam{常考内部/外部碎片区别、首次/最佳/最坏适应算法、伙伴系统合并。}


\concept{动态分区分配}
\begin{points}
  \item 不预先固定分区，进程装入时根据需要动态划分连续内存区。
  \item 优点：比固定分区灵活，内部碎片少。
  \item 缺点：随着分配和回收，会产生外部碎片，需要紧凑/拼接。
\end{points}
\plain{连续分配要求一整块内存，容易产生碎片；分页把连续要求打碎。}
\exam{常考内部/外部碎片区别、首次/最佳/最坏适应算法、伙伴系统合并。}


\concept{内部碎片与外部碎片}
\begin{points}
  \item 内部碎片：已经分配给进程但进程用不完的空间，位于已分配区域内部。
  \item 外部碎片：空闲空间总量足够，但分散在各处，无法形成一个足够大的连续区域。
  \item 固定分区和分页容易有内部碎片；动态分区和分段容易有外部碎片。
  \item 直觉：内部碎片是“房间给你了但你没住满”；外部碎片是“空房很多但都太零散，放不下大团队”。
\end{points}
\plain{连续分配要求一整块内存，容易产生碎片；分页把连续要求打碎。}
\exam{常考内部/外部碎片区别、首次/最佳/最坏适应算法、伙伴系统合并。}


\concept{动态分区分配算法}
\begin{points}
  \item 首次适应 FF：从低地址开始找到第一个足够大的空闲分区。速度较快，但低地址碎片较多。
  \item 循环首次适应 NF：从上次查找结束位置继续查找，分布较均匀。
  \item 最佳适应 BF：选择能满足要求的最小空闲分区，尽量减少剩余空间，但容易产生很多小碎片。
  \item 最坏适应 WF：选择最大的空闲分区，剩余空间仍较大，但可能破坏大空闲区。
\end{points}
\plain{连续分配要求一整块内存，容易产生碎片；分页把连续要求打碎。}
\exam{常考内部/外部碎片区别、首次/最佳/最坏适应算法、伙伴系统合并。}


\concept{伙伴系统}
\begin{points}
  \item 伙伴系统按 $2^k$ 大小管理内存块，申请时分裂大块，释放时若相邻伙伴空闲则合并。
  \item 优点：分裂和合并规则简单，管理效率较高。
  \item 缺点：申请大小不一定正好是 $2^k$，可能产生内部碎片。
\end{points}
\plain{连续分配要求一整块内存，容易产生碎片；分页把连续要求打碎。}
\exam{常考内部/外部碎片区别、首次/最佳/最坏适应算法、伙伴系统合并。}


\section{分页存储管理}
\concept{分页基本思想}
\begin{points}
  \item 将逻辑地址空间划分为大小相等的页 page。
  \item 将物理内存划分为同样大小的页框/页帧 frame。
  \item 以页为单位分配内存，逻辑上连续，物理上可以分散。
  \item 分页取消了作业对物理连续内存的要求，避免外部碎片，但最后一页可能有内部碎片。
\end{points}
\plain{分页就是逻辑上按页切、物理上按页框放，页表负责记录页到页框的映射。}
\exam{常考页号/页内偏移拆分、有效访问时间、多级页表节省页表空间。}


\concept{页表}
\begin{points}
  \item 系统为每个进程建立一张页表，记录逻辑页号到物理块号的映射。
  \item 每个页面对应一个页表项，页表项通常包含物理块号和标志位。
  \item 标志位可能包括存在位、修改位、引用位、保护位等。
  \item 进程切换时，页表基址和长度等信息会被装入页表寄存器。
\end{points}
\plain{分页就是逻辑上按页切、物理上按页框放，页表负责记录页到页框的映射。}
\exam{常考页号/页内偏移拆分、有效访问时间、多级页表节省页表空间。}


\concept{分页地址结构与转换}
\begin{points}
  \item 逻辑地址由页号 $p$ 和页内偏移 $d$ 组成。
  \item 页大小为 $2^n$ 字节时，低 $n$ 位是页内偏移，高位是页号。
  \item 地址转换：用页号查页表得到物理块号 $f$，物理地址为 $f\times$ 页大小 $+d$。
  \item 若页表项无效或越界，则产生地址越界或缺页异常。
\end{points}
\plain{进程不是一直运行，它会在就绪、运行、阻塞等状态之间因调度或等待事件而移动。}
\exam{常考三/五/七状态转换图，尤其运行到阻塞、阻塞到就绪、就绪到运行的触发条件。}

\concept{快表 TLB 与有效访问时间}
\begin{points}
  \item 快表 TLB 是页表项的高速缓存，用于减少访问页表的内存次数。
  \item 无快表时，访问一次数据通常需先访问页表，再访问数据，共两次内存访问。
  \item 有快表时，若命中，直接得到物理块号，只需访问数据；若未命中，需要访问页表再访问数据。
  \item 有效访问时间公式常写为：$EAT=\alpha(t_{TLB}+t_M)+(1-\alpha)(t_{TLB}+2t_M)$，其中 $\alpha$ 为命中率。
\end{points}
\plain{分页就是逻辑上按页切、物理上按页框放，页表负责记录页到页框的映射。}
\exam{常考页号/页内偏移拆分、有效访问时间、多级页表节省页表空间。}


\concept{多级页表}
\begin{points}
  \item 多级页表把页表本身分页，只为实际用到的页表页分配内存。
  \item 解决页表过大、必须连续存放的问题。
  \item 缺点：地址转换需要多次查表，若无 TLB 支持会增加访问开销。
  \item 做题时根据每级索引位数判断每级表项数，根据页大小和表项大小判断一页能装多少表项。
\end{points}
\plain{分页就是逻辑上按页切、物理上按页框放，页表负责记录页到页框的映射。}
\exam{常考页号/页内偏移拆分、有效访问时间、多级页表节省页表空间。}


\section{分段与段页式存储管理}
\concept{分段基本思想}
\begin{points}
  \item 分段按程序的逻辑结构划分地址空间，如代码段、数据段、栈段、共享库段。
  \item 逻辑地址由段号和段内偏移组成。
  \item 段表项包含段基址、段长和保护信息。
  \item 地址转换时先检查段号和段内偏移是否合法，再由段基址加偏移得到物理地址。
\end{points}
\plain{分页按固定大小切，方便管理；分段按程序逻辑模块切，方便共享保护。段页式是先分段再分页。}
\exam{常考分页与分段比较、段号+段内地址合法性、段页式地址结构。}


\concept{分页与分段比较}
\begin{longtable}{p{0.20\textwidth}p{0.36\textwidth}p{0.36\textwidth}}
\toprule
比较项 & 分页 & 分段 \\
\midrule
划分依据 & 系统按固定大小机械划分 & 用户程序逻辑结构划分 \\
单位大小 & 页大小固定 & 段大小可变 \\
地址空间 & 一维地址空间 & 二维地址空间：段号+段内偏移 \\
碎片 & 无外部碎片，有内部碎片 & 无内部碎片或较少，有外部碎片 \\
共享保护 & 不如分段自然 & 以逻辑段为单位，便于共享与保护 \\
\bottomrule
\end{longtable}

\concept{段页式管理}
\begin{points}
  \item 段页式先把程序按逻辑分段，再把每个段分页。
  \item 地址结构通常为段号、段内页号、页内偏移。
  \item 每个段有段表项，段表项指向该段对应的页表。
  \item 优点：兼具分段的逻辑保护共享和分页的离散分配；缺点是地址转换更复杂。
\end{points}
\plain{分页按固定大小切，方便管理；分段按程序逻辑模块切，方便共享保护。段页式是先分段再分页。}
\exam{常考分页与分段比较、段号+段内地址合法性、段页式地址结构。}


\chapter{虚拟存储器}

\section{虚拟存储器概念}
\concept{为什么需要虚拟存储器}
\begin{points}
  \item 传统物理存储管理要求作业运行前全部装入内存，运行期间常驻内存，限制了大作业和多作业运行。
  \item 物理扩充内存会增加成本，因此需要从逻辑上扩充内存。
  \item 虚拟存储器打破“作业必须全部装入内存才能运行”的限制。
\end{points}
\plain{不是资源真的变多，而是 OS 用映射、换入换出、排队等办法，让用户感觉自己有一份独立资源。}
\exam{常考虚拟处理机、虚拟存储器、虚拟设备的例子，以及虚拟性和资源复用的关系。}


\concept{局部性原理}
\begin{points}
  \item 程序执行时，CPU 访问的指令和数据倾向于聚集在较小范围内。
  \item 时间局部性：刚访问过的指令或数据，短时间内可能再次访问。
  \item 空间局部性：访问某地址后，附近地址很可能被访问。
  \item 顺序局部性：程序大多数情况下顺序执行，下一条指令通常紧接当前指令。
  \item 局部性原理是虚拟存储器的理论基础。
\end{points}
\plain{程序一段时间只爱访问附近的一小块代码和数据，所以没必要一开始全装入内存。}
\exam{常考时间/空间局部性、虚拟存储器四特征、容量受地址结构和外存限制。}


\concept{虚拟存储器定义与特征}
\begin{points}
  \item 虚拟存储器是具有请求调入和置换功能、能从逻辑上扩充内存容量的存储器系统。
  \item 它是一种以时间换空间的技术：用缺页中断、磁盘换入换出换取更大的逻辑空间。
  \item 特征：离散性、多次性、对换性、虚拟性。
  \item 容量受地址结构和外存容量限制，不是无限大。
\end{points}
\plain{不是资源真的变多，而是 OS 用映射、换入换出、排队等办法，让用户感觉自己有一份独立资源。}
\exam{常考虚拟处理机、虚拟存储器、虚拟设备的例子，以及虚拟性和资源复用的关系。}


\section{请求分页存储管理}
\concept{请求分页的基本思想}
\begin{points}
  \item 程序运行前只装入部分页面即可启动。
  \item 运行中访问的页面若不在内存，产生缺页中断，由 OS 调入所需页面。
  \item 若内存无空闲页框，需要选择某个页面换出。
  \item 请求分页 = 基本分页 + 请求调页 + 页面置换。
\end{points}
\plain{分页就是逻辑上按页切、物理上按页框放，页表负责记录页到页框的映射。}
\exam{常考页号/页内偏移拆分、有效访问时间、多级页表节省页表空间。}


\concept{基本分页与请求分页区别}
\begin{points}
  \item 基本分页要求进程运行所需页面已在内存中，地址转换时页表项应有效。
  \item 请求分页允许部分页面不在内存，访问未驻留页面时产生缺页中断。
  \item 因此基本分页本身不以缺页中断为常规机制；请求分页才把缺页作为正常运行过程的一部分。
\end{points}
\plain{分页就是逻辑上按页切、物理上按页框放，页表负责记录页到页框的映射。}
\exam{常考页号/页内偏移拆分、有效访问时间、多级页表节省页表空间。}


\concept{页表项扩展}
\begin{points}
  \item 存在位/有效位：表示该页是否在内存。
  \item 修改位/脏位：表示页面是否被写过，换出时若为脏页需写回外存。
  \item 引用位/访问位：表示页面近期是否被访问，用于页面置换算法。
  \item 外存地址：指出页面在外存中的位置。
\end{points}
\plain{分页就是逻辑上按页切、物理上按页框放，页表负责记录页到页框的映射。}
\exam{常考页号/页内偏移拆分、有效访问时间、多级页表节省页表空间。}


\concept{缺页中断处理过程}
\begin{points}
  \item CPU 访问页面，发现页表项存在位为 0，产生缺页中断。
  \item 操作系统检查访问是否合法；若非法，终止进程。
  \item 若合法，查找空闲页框；若无空闲页框，执行页面置换。
  \item 若被换出页被修改过，需要写回外存。
  \item 从外存读入所需页面，更新页表和快表，重新执行引起缺页的指令。
\end{points}
\plain{中断就是硬件或软件打断 CPU 当前执行，让 OS 先处理更紧急或必须处理的事件。}
\exam{常考中断处理步骤、中断与异常/系统调用区别、为什么中断能提高 CPU 与 I/O 并行度。}


\section{页面置换算法}
\concept{最佳置换 OPT}
\begin{points}
  \item OPT 选择未来最长时间不会被访问的页面淘汰。
  \item 它缺页率最低，但需要预知未来，实际无法实现。
  \item 常作为评价其他算法的理论下界。
\end{points}
\plain{内存满了要踢页；OPT 看未来，FIFO 看入队时间，LRU 看最近使用，CLOCK 用引用位近似 LRU。}
\exam{常考给页面引用串和页框数计算缺页次数，FIFO 可能有 Belady 异常。}


\concept{先进先出 FIFO}
\begin{points}
  \item FIFO 淘汰最早进入内存的页面。
  \item 实现简单，只需维护队列。
  \item 缺点是不考虑页面使用频率，可能淘汰仍常用页面。
  \item FIFO 可能出现 Belady 异常：页框数增加，缺页次数反而增加。
\end{points}
\plain{内存满了要踢页；OPT 看未来，FIFO 看入队时间，LRU 看最近使用，CLOCK 用引用位近似 LRU。}
\exam{常考给页面引用串和页框数计算缺页次数，FIFO 可能有 Belady 异常。}


\concept{最近最久未使用 LRU}
\begin{points}
  \item LRU 淘汰最近最长时间没有被访问的页面。
  \item 基于局部性原理：最近很久未用的页面，未来短期内也可能不用。
  \item 实现可用计数器、栈或近似算法，但精确 LRU 开销较高。
  \item LRU 不会出现 Belady 异常。
\end{points}
\plain{内存满了要踢页；OPT 看未来，FIFO 看入队时间，LRU 看最近使用，CLOCK 用引用位近似 LRU。}
\exam{常考给页面引用串和页框数计算缺页次数，FIFO 可能有 Belady 异常。}


\concept{时钟 CLOCK/二次机会算法}
\begin{points}
  \item Clock 用访问位近似 LRU，把页面组成循环队列。
  \item 指针扫描页面：若访问位为 0，淘汰；若为 1，清 0 并给第二次机会。
  \item 改进型 Clock 还考虑修改位，优先淘汰未访问且未修改的页面。
\end{points}
\plain{内存满了要踢页；OPT 看未来，FIFO 看入队时间，LRU 看最近使用，CLOCK 用引用位近似 LRU。}
\exam{常考给页面引用串和页框数计算缺页次数，FIFO 可能有 Belady 异常。}


\section{页面分配与置换策略}
\concept{驻留集与页面分配}
\begin{points}
  \item 驻留集指某进程当前驻留内存的页面集合。
  \item 固定分配：进程运行期间分得的页框数固定。
  \item 可变分配：根据缺页率、工作集等动态调整页框数。
\end{points}
\plain{把这个知识点放回本章主线理解：它回答的是 OS 为了管理资源、隐藏硬件、保证并发正确性而采用的某个对象、规则或步骤。}
\exam{常考选择/填空/简答，让你判断概念归属、比较相近概念，或按“定义-作用-机制-易错点”展开。}


\concept{局部置换与全局置换}
\begin{points}
  \item 局部置换：缺页时只从本进程驻留集中选择页面淘汰。
  \item 全局置换：可从全系统所有可换页面中选择淘汰，进程页框数会动态变化。
  \item 固定分配只能配合局部置换；可变分配可配合局部置换或全局置换。
\end{points}
\plain{内存满了要踢页；OPT 看未来，FIFO 看入队时间，LRU 看最近使用，CLOCK 用引用位近似 LRU。}
\exam{常考给页面引用串和页框数计算缺页次数，FIFO 可能有 Belady 异常。}


\concept{不能组合的策略}
\begin{points}
  \item 固定分配 + 全局置换通常不能组合。
  \item 原因：固定分配要求每个进程页框数不变，而全局置换可能把其他进程页框拿走或让本进程获得额外页框，改变驻留集大小。
\end{points}
\plain{把这个知识点放回本章主线理解：它回答的是 OS 为了管理资源、隐藏硬件、保证并发正确性而采用的某个对象、规则或步骤。}
\exam{常考选择/填空/简答，让你判断概念归属、比较相近概念，或按“定义-作用-机制-易错点”展开。}


\section{工作集、抖动与缺页率}
\concept{工作集模型}
\begin{points}
  \item 工作集是在某段时间窗口内进程实际访问过的页面集合。
  \item 工作集反映进程当前局部性需求，系统应尽量给进程分配足够页框容纳其工作集。
  \item 若分配页框少于工作集需要，缺页率会显著上升。
\end{points}
\plain{给进程的页框太少，它会不停缺页，把 CPU 时间浪费在换页上，这就是抖动。}
\exam{常考抖动原因和解决：降低多道程序度、增加页框、按工作集分配。}


\concept{抖动 thrashing}
\begin{points}
  \item 抖动指系统花费大量时间进行页面换入换出，真正执行用户程序的时间很少。
  \item 现象：CPU 利用率低，磁盘交换区利用率很高，缺页频繁。
  \item 原因：多道程序度过高或每个进程分得页框不足，工作集无法驻留内存。
  \item 解决：降低多道程序度、增加页框、采用工作集策略、改善置换算法。
\end{points}
\plain{给进程的页框太少，它会不停缺页，把 CPU 时间浪费在换页上，这就是抖动。}
\exam{常考抖动原因和解决：降低多道程序度、增加页框、按工作集分配。}


\concept{影响缺页率的因素}
\begin{points}
  \item 页面大小：过小会增加页表和 I/O 次数，过大会增加内部碎片和无用调入。
  \item 分配页框数：一般页框越多缺页率越低，但 FIFO 可能有 Belady 异常。
  \item 页面置换算法：OPT 最优，LRU 通常较好，FIFO 简单但可能异常。
  \item 程序局部性：局部性越好，缺页率越低。
  \item 多道程序度：过高会导致每个进程页框不足，引发抖动。
\end{points}
\plain{页不在内存时才调入；这会产生缺页中断，由 OS 找页、分配页框、必要时置换。}
\exam{常考基本分页与请求分页区别、缺页中断流程、页表项存在位/修改位/引用位。}


\chapter{文件系统}

\section{文件基本概念}
\concept{文件系统的组成}
\begin{points}
  \item 文件系统是操作系统中与管理文件有关的软件和数据的集合。
  \item 组成包括：被管理的文件集合、管理文件所需的数据结构、相应管理软件。
  \item 管理软件负责外存空间分配回收、目录管理、文件读写、共享与保护等。
\end{points}
\plain{文件是用户看到的逻辑外存单位；用户不用管扇区，只管名字、属性和读写操作。}
\exam{常考文件属性、文件操作、open 为什么能减少目录搜索次数。}


\concept{文件定义}
\begin{points}
  \item 文件是记录在外部存储上，具有名字，在逻辑上有完整意义的信息项序列。
  \item 信息项是构成文件内容的基本单位，可以是字节或记录。
  \item 从用户角度看，文件是逻辑外存的最小分配单元。
  \item 文件内容的意义由文件创建者和使用者解释。
\end{points}
\plain{文件是用户看到的逻辑外存单位；用户不用管扇区，只管名字、属性和读写操作。}
\exam{常考文件属性、文件操作、open 为什么能减少目录搜索次数。}


\concept{文件属性}
\begin{points}
  \item 名称：便于用户识别的符号文件名。
  \item 标识符：系统内部唯一标识，常为数字形式。
  \item 类型：普通文件、目录文件、设备文件等。
  \item 位置：文件在设备上的位置指针。
  \item 大小：当前大小和最大允许大小。
  \item 保护：访问控制信息，规定谁能读、写、执行。
  \item 时间、日期、用户标识：创建、修改、访问等元数据。
\end{points}
\plain{文件是用户看到的逻辑外存单位；用户不用管扇区，只管名字、属性和读写操作。}
\exam{常考文件属性、文件操作、open 为什么能减少目录搜索次数。}


\concept{文件操作}
\begin{points}
  \item Create：创建文件，为文件分配目录项和必要空间。
  \item Write：根据写指针向文件写入数据。
  \item Read：根据读指针从文件读取数据。
  \item Seek：移动文件当前读写位置。
  \item Delete：删除文件，释放空间，删除目录项。
  \item Truncate：删除文件内容但保留属性。
  \item Open/Close：打开文件时把文件控制信息装入内存打开文件表，关闭时释放相关表项。
\end{points}
\plain{文件是用户看到的逻辑外存单位；用户不用管扇区，只管名字、属性和读写操作。}
\exam{常考文件属性、文件操作、open 为什么能减少目录搜索次数。}


\section{文件访问方法与逻辑结构}
\concept{顺序访问}
\begin{points}
  \item 顺序访问按文件中信息项的先后顺序读写。
  \item 适合文本流、日志、磁带等场景。
  \item 若要访问第 $n$ 个记录，通常需要从前面依次读到它。
\end{points}
\plain{文件怎么读有两层：逻辑上记录如何组织，物理上块如何放在磁盘。}
\exam{常考顺序/直接访问，顺序文件、索引文件、索引顺序文件、散列文件适用场景。}


\concept{直接访问/随机访问}
\begin{points}
  \item 直接访问允许按记录号或块号直接定位到某个位置。
  \item 磁盘支持随机访问，因此文件系统可实现直接访问。
  \item 适合数据库、索引文件、需要频繁定位的数据文件。
\end{points}
\plain{文件怎么读有两层：逻辑上记录如何组织，物理上块如何放在磁盘。}
\exam{常考顺序/直接访问，顺序文件、索引文件、索引顺序文件、散列文件适用场景。}


\concept{文件逻辑结构}
\begin{points}
  \item 无结构文件/流式文件：文件被看作字节序列，如普通文本和二进制文件。
  \item 有结构文件/记录式文件：文件由一组记录组成，每条记录有内部字段结构。
  \item 顺序文件：记录按某种顺序排列，适合顺序处理。
  \item 索引文件：建立索引表，从键值映射到记录位置，提高随机查找效率。
  \item 索引顺序文件：主文件按键有序，同时建立分组索引，兼顾顺序处理和快速定位。
  \item 散列文件：通过散列函数直接计算记录位置，适合等值查找。
\end{points}
\plain{文件怎么读有两层：逻辑上记录如何组织，物理上块如何放在磁盘。}
\exam{常考顺序/直接访问，顺序文件、索引文件、索引顺序文件、散列文件适用场景。}


\section{目录结构}
\concept{目录与目录项}
\begin{points}
  \item 目录用于组织系统中的文件并提供文件信息。
  \item 目录项通常包含文件名和文件控制块指针或 inode 号。
  \item 目录检索就是根据路径名逐级查找目录项，找到目标文件控制信息。
\end{points}
\plain{目录就是文件名到文件控制信息的映射表，路径就是沿目录树找到文件。}
\exam{常考一级/两级/树形目录优缺点、平均查找磁盘访问次数。}


\concept{一级目录}
\begin{points}
  \item 系统只有一个目录，所有文件都放在其中。
  \item 优点：实现简单。
  \item 缺点：文件多时查找慢，不能解决重名问题，不便于多用户管理。
\end{points}
\plain{目录就是文件名到文件控制信息的映射表，路径就是沿目录树找到文件。}
\exam{常考一级/两级/树形目录优缺点、平均查找磁盘访问次数。}


\concept{两级目录}
\begin{points}
  \item 每个用户有自己的用户文件目录 UFD，系统有主文件目录 MFD。
  \item 解决不同用户之间文件重名问题。
  \item 用户内部文件多时仍不够灵活。
\end{points}
\plain{目录就是文件名到文件控制信息的映射表，路径就是沿目录树找到文件。}
\exam{常考一级/两级/树形目录优缺点、平均查找磁盘访问次数。}


\concept{树形目录与路径}
\begin{points}
  \item 树形目录把目录组织成层次结构，便于分类管理。
  \item 绝对路径从根目录开始，如 \code{/home/a.txt}。
  \item 相对路径从当前目录开始。
  \item 优点：层次清晰、重名控制灵活、便于权限管理。
\end{points}
\plain{目录就是文件名到文件控制信息的映射表，路径就是沿目录树找到文件。}
\exam{常考一级/两级/树形目录优缺点、平均查找磁盘访问次数。}


\concept{无环图目录与共享}
\begin{points}
  \item 无环图目录允许多个目录项共享同一个文件或子目录，但不允许形成环。
  \item 共享文件需要引用计数，记录有多少目录项指向该文件。
  \item 引用计数为 0 时，文件数据和控制信息才能真正回收。
\end{points}
\plain{多个进程都想用同一批资源，所以必须规定能不能一起用、谁先用、用完怎么释放。}
\exam{常考互斥共享/同时共享举例，或者问并发和共享为什么是 OS 最基本特征。}


\section{文件系统安装、共享与保护}
\concept{文件系统安装 mount}
\begin{points}
  \item 文件系统必须挂载到某个目录位置后，用户才能通过统一目录树访问它。
  \item 挂载点是原有目录树中的一个目录。
  \item 挂载后，该目录下呈现被挂载文件系统的根内容。
\end{points}
\plain{文件是用户看到的逻辑外存单位；用户不用管扇区，只管名字、属性和读写操作。}
\exam{常考文件属性、文件操作、open 为什么能减少目录搜索次数。}


\concept{硬链接与符号链接}
\begin{points}
  \item 硬链接是多个目录项指向同一个 inode/文件控制块，本质上是同一个文件有多个名字。
  \item 删除一个硬链接只减少引用计数；只有引用计数为 0 时文件才真正删除。
  \item 符号链接是一个特殊文件，内容是目标文件路径。
  \item 符号链接可以跨文件系统，也可以指向目录；但目标删除后可能变成悬空链接。
\end{points}
\plain{硬链接共享同一个 inode，符号链接是另一个文件，内容是目标路径。打开文件表保存运行时读写状态。}
\exam{常考硬/软链接区别、引用计数、系统级/进程级打开文件表。}


\concept{文件保护}
\begin{points}
  \item 文件保护用于控制用户对文件的读、写、执行、删除等权限。
  \item 常见方法包括访问控制表 ACL、权限位、口令保护、能力表等。
  \item Unix 风格权限常分为所有者、同组用户、其他用户三类，每类有读写执行位。
\end{points}
\plain{把这个知识点放回本章主线理解：它回答的是 OS 为了管理资源、隐藏硬件、保证并发正确性而采用的某个对象、规则或步骤。}
\exam{常考选择/填空/简答，让你判断概念归属、比较相近概念，或按“定义-作用-机制-易错点”展开。}


\concept{打开文件表}
\begin{points}
  \item 打开文件后，系统把文件控制信息缓存在内存打开文件表中，避免每次读写都重新搜索目录。
  \item 系统级打开文件表记录文件位置、引用计数、磁盘位置等全局信息。
  \item 进程级打开文件表记录进程自己的文件描述符、访问模式、当前读写指针等。
  \item 若父子进程共享同一个打开文件表项，可能共享文件偏移；若分别打开同一文件，则偏移独立。
\end{points}
\plain{硬链接共享同一个 inode，符号链接是另一个文件，内容是目标路径。打开文件表保存运行时读写状态。}
\exam{常考硬/软链接区别、引用计数、系统级/进程级打开文件表。}


\chapter{文件系统实现}

\section{文件系统结构}
\concept{用户视图与系统视图}
\begin{points}
  \item 用户视图中，文件是带名字的逻辑数据集合，通过系统调用访问。
  \item 系统视图中，文件是磁盘数据块的集合，文件系统负责从逻辑文件偏移映射到物理磁盘块。
  \item 文件系统基本操作单位是数据块 block，而磁盘物理基本单位是扇区 sector。
  \item 操作系统通常一次读写一个块，块可由多个连续扇区组成。
\end{points}
\plain{把这个知识点放回本章主线理解：它回答的是 OS 为了管理资源、隐藏硬件、保证并发正确性而采用的某个对象、规则或步骤。}
\exam{常考选择/填空/简答，让你判断概念归属、比较相近概念，或按“定义-作用-机制-易错点”展开。}


\concept{分层结构}
\begin{points}
  \item 逻辑文件系统：处理文件名、目录、权限、FCB/inode 等。
  \item 文件组织模块：完成逻辑块到物理块的映射。
  \item 基本文件系统：向设备驱动发出读写块请求。
  \item I/O 控制层：设备驱动和中断处理程序控制硬件。
\end{points}
\plain{把这个知识点放回本章主线理解：它回答的是 OS 为了管理资源、隐藏硬件、保证并发正确性而采用的某个对象、规则或步骤。}
\exam{常考选择/填空/简答，让你判断概念归属、比较相近概念，或按“定义-作用-机制-易错点”展开。}


\concept{FCB 与 inode}
\begin{points}
  \item FCB 文件控制块保存文件属性和存储位置信息。
  \item inode 是类 Unix 系统中的索引节点，保存文件类型、权限、大小、时间、数据块地址等元数据。
  \item 目录项通常保存文件名和 inode 号，文件名不一定存放在 inode 内。
  \item inode 与文件物理分配方式有关，因为 inode 中需要保存指向数据块或索引块的指针。
\end{points}
\plain{FCB/inode 保存文件属性和物理位置，是文件系统找到文件内容的核心索引。}
\exam{常考目录项是否要放全部 FCB、inode 与文件名分离的好处。}


\section{目录实现}
\concept{线性列表目录}
\begin{points}
  \item 目录项按线性表保存，查找时顺序扫描。
  \item 优点：实现简单。
  \item 缺点：目录项多时查找慢，平均查找约需扫描一半目录项。
  \item 目录项访问次数题常用“目录项数/每块目录项数/平均扫描块数”计算。
\end{points}
\plain{目录就是文件名到文件控制信息的映射表，路径就是沿目录树找到文件。}
\exam{常考一级/两级/树形目录优缺点、平均查找磁盘访问次数。}


\concept{哈希目录}
\begin{points}
  \item 根据文件名计算哈希值，快速定位目录项。
  \item 优点：查找速度快。
  \item 缺点：需要处理哈希冲突，目录扩展和哈希表维护更复杂。
\end{points}
\plain{目录就是文件名到文件控制信息的映射表，路径就是沿目录树找到文件。}
\exam{常考一级/两级/树形目录优缺点、平均查找磁盘访问次数。}


\section{文件存储空间分配方法}
\concept{连续分配}
\begin{points}
  \item 文件占用磁盘上一组连续块，目录项只需记录起始块号和长度。
  \item 优点：顺序访问速度快，随机访问也容易，寻道少。
  \item 缺点：需要预先知道文件大小，文件增长困难，会产生外部碎片。
  \item 适合大小较稳定、顺序访问多的文件。
\end{points}
\plain{连续分配要求一整块内存，容易产生碎片；分页把连续要求打碎。}
\exam{常考内部/外部碎片区别、首次/最佳/最坏适应算法、伙伴系统合并。}


\concept{链式分配}
\begin{points}
  \item 文件由多个磁盘块组成链表，每个块包含指向下一个块的指针。
  \item 优点：无外部碎片，文件容易增长。
  \item 缺点：随机访问差，访问第 $k$ 个块必须从第一个块顺链找；指针损坏会影响后续数据。
  \item FAT 是链式分配的改进，把链接指针集中存放在文件分配表中。
\end{points}
\plain{每个块指向下一个块，不怕文件增长，但找第 n 块必须顺着链走。}
\exam{常考链式分配不适合随机访问，FAT 如何把指针集中到表中改善。}


\concept{索引分配}
\begin{points}
  \item 为每个文件建立索引块，索引块中保存该文件所有数据块号。
  \item 目录项保存索引块地址。
  \item 优点：支持直接访问，无外部碎片，文件可离散存放。
  \item 缺点：索引块本身占空间；小文件可能浪费索引块，大文件需要多级索引。
\end{points}
\plain{把文件所有数据块号集中放在索引块里，既支持离散分配，又方便随机访问。}
\exam{常考单级/多级索引最大文件长度计算：每个索引块能放多少块号，再乘数据块大小。}


\concept{多级索引与最大文件长度}
\begin{points}
  \item 若块大小为 $B$ 字节，每个块号占 $p$ 字节，则一个索引块可保存 $B/p$ 个块号。
  \item 一级索引最大可索引 $(B/p)$ 个数据块，最大文件大小为 $(B/p)\times B$。
  \item 二级索引最大可索引 $(B/p)^2$ 个数据块，最大文件大小为 $(B/p)^2\times B$。
  \item 三级索引同理为 $(B/p)^3\times B$。
\end{points}
\plain{把文件所有数据块号集中放在索引块里，既支持离散分配，又方便随机访问。}
\exam{常考单级/多级索引最大文件长度计算：每个索引块能放多少块号，再乘数据块大小。}

\section{空闲空间管理}
\concept{位示图/位图}
\begin{points}
  \item 用一个二进制位表示一个磁盘块是否空闲。
  \item 优点：查找连续空闲块方便，结构紧凑。
  \item 缺点：磁盘很大时位图本身也较大，需要常驻或缓存。
\end{points}
\plain{文件系统要知道哪些盘块空着，常用位图、链表、成组链接等方法管理。}
\exam{常考位示图位号与盘块号换算、空闲链表/成组链接特点。}


\concept{空闲链表}
\begin{points}
  \item 把所有空闲块链接成链表。
  \item 优点：释放和分配单个块较简单。
  \item 缺点：查找连续空闲块效率低，链指针损坏可能影响管理。
\end{points}
\plain{文件系统要知道哪些盘块空着，常用位图、链表、成组链接等方法管理。}
\exam{常考位示图位号与盘块号换算、空闲链表/成组链接特点。}


\concept{成组链接法}
\begin{points}
  \item 把空闲块分组，每组第一块记录下一组空闲块号。
  \item 兼顾链表和批量管理，是 Unix 早期常用方法。
  \item 适合大量空闲块的快速分配和回收。
\end{points}
\plain{硬链接共享同一个 inode，符号链接是另一个文件，内容是目标路径。打开文件表保存运行时读写状态。}
\exam{常考硬/软链接区别、引用计数、系统级/进程级打开文件表。}


\concept{计数法}
\begin{points}
  \item 用“起始空闲块号 + 连续空闲块数”表示一段连续空闲区。
  \item 适合磁盘空闲块成片连续的情况。
  \item 比逐块链表更紧凑，但碎片多时记录数量增加。
\end{points}
\plain{把这个知识点放回本章主线理解：它回答的是 OS 为了管理资源、隐藏硬件、保证并发正确性而采用的某个对象、规则或步骤。}
\exam{常考选择/填空/简答，让你判断概念归属、比较相近概念，或按“定义-作用-机制-易错点”展开。}


\section{文件共享实现}
\concept{基于索引节点的共享}
\begin{points}
  \item 多个目录项指向同一个 inode，实现硬链接共享。
  \item inode 中维护引用计数，记录有多少目录项引用它。
  \item 删除目录项只减少引用计数，不一定删除文件内容。
\end{points}
\plain{多个进程都想用同一批资源，所以必须规定能不能一起用、谁先用、用完怎么释放。}
\exam{常考互斥共享/同时共享举例，或者问并发和共享为什么是 OS 最基本特征。}


\concept{基于符号链接的共享}
\begin{points}
  \item 符号链接文件保存目标路径名。
  \item 访问符号链接时，系统根据路径再去查找目标文件。
  \item 优点是灵活、可跨文件系统；缺点是访问多一次路径解析，目标删除后链接失效。
\end{points}
\plain{多个进程都想用同一批资源，所以必须规定能不能一起用、谁先用、用完怎么释放。}
\exam{常考互斥共享/同时共享举例，或者问并发和共享为什么是 OS 最基本特征。}


\chapter{大容量存储系统}

\section{磁盘结构与访问时间}
\concept{磁盘结构}
\begin{points}
  \item 磁盘由一个或多个盘片组成，盘片表面覆盖磁性材料。
  \item 每个盘面对应一个磁头，磁臂沿半径方向移动。
  \item 磁道是盘面上的同心圆，扇区是磁道上的固定大小区域。
  \item 柱面是所有盘面上半径相同的磁道集合。
  \item 一个物理块可用柱面号、磁头号、扇区号表示。
\end{points}
\plain{磁盘访问慢主要慢在机械移动：寻道和旋转。调度算法本质是少移动磁臂。}
\exam{常考按磁道请求序列计算移动磁道数，区分 SCAN/C-SCAN/LOOK/C-LOOK。}


\concept{磁盘访问时间}
\begin{points}
  \item 寻道时间：磁臂移动到目标柱面/磁道所需时间。
  \item 旋转延迟：目标扇区旋转到磁头下方所需时间，平均为旋转一周时间的一半。
  \item 传输时间：数据从磁盘读出或写入磁盘的时间。
  \item 总访问时间约为：寻道时间 + 旋转延迟 + 传输时间。
  \item 机械磁盘优化重点通常是减少寻道时间。
\end{points}
\plain{文件怎么读有两层：逻辑上记录如何组织，物理上块如何放在磁盘。}
\exam{常考顺序/直接访问，顺序文件、索引文件、索引顺序文件、散列文件适用场景。}


\section{磁盘调度算法}
\concept{FCFS}
\begin{points}
  \item 按磁盘请求到达顺序服务。
  \item 优点：公平、简单。
  \item 缺点：磁头来回移动可能很大，平均寻道距离长。
\end{points}
\plain{谁先来谁先上，最朴素但容易让短作业排在长作业后面。}
\exam{常考护航效应和平均等待时间计算。}


\concept{SSTF 最短寻道时间优先}
\begin{points}
  \item 每次选择距离当前磁头位置最近的请求。
  \item 优点：平均寻道距离通常较短。
  \item 缺点：远处请求可能长期等待，产生饥饿。
\end{points}
\plain{磁盘访问慢主要慢在机械移动：寻道和旋转。调度算法本质是少移动磁臂。}
\exam{常考按磁道请求序列计算移动磁道数，区分 SCAN/C-SCAN/LOOK/C-LOOK。}


\concept{SCAN 电梯算法}
\begin{points}
  \item 磁头像电梯一样沿一个方向移动，服务沿途请求，到达端点后反向。
  \item 优点：比 FCFS 稳定，避免频繁来回。
  \item 缺点：靠近两端的请求等待时间可能不同。
\end{points}
\plain{磁盘访问慢主要慢在机械移动：寻道和旋转。调度算法本质是少移动磁臂。}
\exam{常考按磁道请求序列计算移动磁道数，区分 SCAN/C-SCAN/LOOK/C-LOOK。}


\concept{C-SCAN 循环扫描}
\begin{points}
  \item 磁头只沿一个方向服务，到达一端后快速返回另一端，不在返回途中服务。
  \item 提供更均匀的等待时间。
  \item 适合请求分布较均匀的系统。
\end{points}
\plain{磁盘访问慢主要慢在机械移动：寻道和旋转。调度算法本质是少移动磁臂。}
\exam{常考按磁道请求序列计算移动磁道数，区分 SCAN/C-SCAN/LOOK/C-LOOK。}


\concept{LOOK 与 C-LOOK}
\begin{points}
  \item LOOK 不一定走到磁盘端点，只走到当前方向最远请求处就反向。
  \item C-LOOK 只在有请求的范围内循环扫描，从一端最远请求跳到另一端最远请求。
  \item 它们是 SCAN/C-SCAN 的实际优化版本，减少无意义移动。
\end{points}
\plain{磁盘访问慢主要慢在机械移动：寻道和旋转。调度算法本质是少移动磁臂。}
\exam{常考按磁道请求序列计算移动磁道数，区分 SCAN/C-SCAN/LOOK/C-LOOK。}


\section{磁盘管理与容错}
\concept{磁盘格式化}
\begin{points}
  \item 低级格式化把磁盘划分为扇区，为每个扇区建立控制信息。
  \item 分区把磁盘划分为一个或多个逻辑区域。
  \item 高级格式化/逻辑格式化在分区上建立文件系统数据结构，如空闲空间表、根目录、inode 区等。
\end{points}
\plain{磁盘访问慢主要慢在机械移动：寻道和旋转。调度算法本质是少移动磁臂。}
\exam{常考按磁道请求序列计算移动磁道数，区分 SCAN/C-SCAN/LOOK/C-LOOK。}


\concept{引导块与坏块}
\begin{points}
  \item 引导块保存启动操作系统所需的引导代码或引导信息。
  \item 坏块是无法可靠读写的磁盘块。
  \item 坏块管理可通过扇区备用、坏块链表、控制器重映射等方式处理。
\end{points}
\plain{磁盘管理不仅是读写，还包括初始化、启动、坏块处理和容错。}
\exam{常考低级/高级格式化、引导块作用、RAID 提高可靠性/性能的方式。}


\concept{交换空间管理}
\begin{points}
  \item 交换空间用于保存被换出的进程映像或页面。
  \item 交换空间可作为专门分区，也可作为文件系统中的交换文件。
  \item 专门分区管理简单且效率高；交换文件更灵活但可能受文件系统开销影响。
\end{points}
\plain{把这个知识点放回本章主线理解：它回答的是 OS 为了管理资源、隐藏硬件、保证并发正确性而采用的某个对象、规则或步骤。}
\exam{常考选择/填空/简答，让你判断概念归属、比较相近概念，或按“定义-作用-机制-易错点”展开。}


\concept{RAID 磁盘容错}
\begin{points}
  \item RAID 通过磁盘冗余提高性能和可靠性。
  \item RAID 0 条带化提高性能但无冗余。
  \item RAID 1 镜像保存完整副本，可靠性高但空间开销大。
  \item RAID 5 使用分布式奇偶校验，兼顾容量和容错。
  \item RAID 6 可容忍两块磁盘故障，可靠性更高。
\end{points}
\plain{磁盘访问慢主要慢在机械移动：寻道和旋转。调度算法本质是少移动磁臂。}
\exam{常考按磁道请求序列计算移动磁道数，区分 SCAN/C-SCAN/LOOK/C-LOOK。}


\chapter{设备管理}

\section{I/O 系统与设备分类}
\concept{设备管理的对象与任务}
\begin{points}
  \item 除 CPU 和内存外，大部分硬件设备称为外部设备。
  \item I/O 系统包括 I/O 设备、接口线路、控制部件、通道和管理软件。
  \item I/O 操作是主存与外围设备介质之间的信息传送。
  \item 设备管理任务：响应 I/O 请求、分配 I/O 设备、管理设备可靠性与安全性、提供友好访问接口。
\end{points}
\plain{I/O 部分就是解决 CPU 快、设备慢、设备复杂且共享的问题。DMA/通道减少 CPU 搬运，缓冲平滑速度差，SPOOLing 把独占设备虚拟成共享设备。}
\exam{常考设备分类、四种 I/O 控制方式、DMA 与中断/通道区别、缓冲和 SPOOLing 的作用。}


\concept{设备管理四个功能}
\begin{points}
  \item 设备分配：根据用户 I/O 请求分配设备、控制器和通道。
  \item 设备处理：启动设备，并在 I/O 完成时处理中断。
  \item 缓冲管理：缓和 CPU 与 I/O 速度不匹配，管理缓冲区分配与释放。
  \item 设备独立性：用户编程使用逻辑设备名，不依赖具体物理设备。
\end{points}
\plain{I/O 部分就是解决 CPU 快、设备慢、设备复杂且共享的问题。DMA/通道减少 CPU 搬运，缓冲平滑速度差，SPOOLing 把独占设备虚拟成共享设备。}
\exam{常考设备分类、四种 I/O 控制方式、DMA 与中断/通道区别、缓冲和 SPOOLing 的作用。}


\concept{按使用特性分类}
\begin{points}
  \item 存储设备：保存信息，如磁盘、磁带、SSD。
  \item 传输设备：传输数据，如网卡、调制解调器。
  \item 人机交互设备：输入或输出用户信息，如键盘、显示器、打印机。
\end{points}
\plain{把这个知识点放回本章主线理解：它回答的是 OS 为了管理资源、隐藏硬件、保证并发正确性而采用的某个对象、规则或步骤。}
\exam{常考选择/填空/简答，让你判断概念归属、比较相近概念，或按“定义-作用-机制-易错点”展开。}


\concept{按共享属性分类}
\begin{points}
  \item 独占设备：一段时间内只允许一个进程使用，如传统打印机。
  \item 共享设备：允许多个进程交替或并发使用，如磁盘。
  \item 虚拟设备：通过虚拟技术把独占设备改造成多个逻辑设备，如 SPOOLing 打印机。
\end{points}
\plain{多个进程都想用同一批资源，所以必须规定能不能一起用、谁先用、用完怎么释放。}
\exam{常考互斥共享/同时共享举例，或者问并发和共享为什么是 OS 最基本特征。}


\concept{按信息交换单位分类}
\begin{points}
  \item 块设备以固定大小数据块为单位传输，如磁盘、磁带。
  \item 字符设备以字符流为单位传输，如键盘、鼠标、串口、打印机。
  \item 选择题中，磁盘通常是块设备；键盘、打印机、显示器通常是字符设备。
\end{points}
\plain{把这个知识点放回本章主线理解：它回答的是 OS 为了管理资源、隐藏硬件、保证并发正确性而采用的某个对象、规则或步骤。}
\exam{常考选择/填空/简答，让你判断概念归属、比较相近概念，或按“定义-作用-机制-易错点”展开。}


\section{I/O 硬件}
\concept{设备控制器与寄存器}
\begin{points}
  \item 设备控制器是 CPU 与外设之间的电子部件，负责控制设备操作。
  \item 控制器通常包含数据寄存器、状态寄存器、控制寄存器。
  \item CPU 通过读写控制器寄存器与设备交互。
  \item 设备完成操作或发生错误时，通过中断通知 CPU。
\end{points}
\plain{硬件基础不用背太深，重点是知道哪些硬件机制支撑 OS：寄存器保存现场，PSW 保存状态，CPU 模式限制权限。}
\exam{常考 CPU 工作流程、单/多处理器、寄存器类别、PSW 中断位/状态位。}


\concept{I/O 地址空间}
\begin{points}
  \item 端口映射 I/O：I/O 设备有独立地址空间，CPU 使用专门 I/O 指令访问端口。
  \item 内存映射 I/O：设备寄存器映射到内存地址空间，CPU 用普通访存指令访问。
  \item 二者都需要保护，普通用户程序不能随意访问设备寄存器。
\end{points}
\plain{程序看到的是逻辑地址，内存条上真正用的是物理地址，中间靠装载器/MMU 转换。}
\exam{常考逻辑地址到物理地址转换、静态/动态重定位、越界检查。}


\concept{通道}
\begin{points}
  \item 通道是一种专门负责 I/O 控制的处理部件，比普通设备控制器更强。
  \item 通道能执行通道程序，组织多个 I/O 操作，减轻 CPU 负担。
  \item 处理机和通道可以并行工作，因此通道提高了大型系统 I/O 并行能力。
\end{points}
\plain{I/O 部分就是解决 CPU 快、设备慢、设备复杂且共享的问题。DMA/通道减少 CPU 搬运，缓冲平滑速度差，SPOOLing 把独占设备虚拟成共享设备。}
\exam{常考设备分类、四种 I/O 控制方式、DMA 与中断/通道区别、缓冲和 SPOOLing 的作用。}


\section{I/O 控制方式}
\concept{程序直接控制方式}
\begin{points}
  \item CPU 反复查询设备状态寄存器，直到设备就绪，再进行数据传送。
  \item 优点：实现简单。
  \item 缺点：CPU 忙等，浪费处理机时间。
  \item 又称轮询方式，适合简单或低速设备。
\end{points}
\plain{把这个知识点放回本章主线理解：它回答的是 OS 为了管理资源、隐藏硬件、保证并发正确性而采用的某个对象、规则或步骤。}
\exam{常考选择/填空/简答，让你判断概念归属、比较相近概念，或按“定义-作用-机制-易错点”展开。}


\concept{中断驱动方式}
\begin{points}
  \item CPU 启动 I/O 后转去执行其他任务，设备完成后发中断通知 CPU。
  \item 优点：减少 CPU 忙等，提高并发性。
  \item 缺点：每次传送仍需 CPU 参与，频繁中断会增加开销。
\end{points}
\plain{中断就是硬件或软件打断 CPU 当前执行，让 OS 先处理更紧急或必须处理的事件。}
\exam{常考中断处理步骤、中断与异常/系统调用区别、为什么中断能提高 CPU 与 I/O 并行度。}


\concept{DMA 方式}
\begin{points}
  \item DMA 允许设备控制器在主存和设备之间直接传送一批数据，无需 CPU 逐字节搬运。
  \item CPU 只负责初始化 DMA 控制器，设置内存地址、传输长度、方向等。
  \item 传输完成后 DMA 控制器发中断通知 CPU。
  \item 优点：大幅减少 CPU 参与，适合磁盘、网卡等高速块设备。
\end{points}
\plain{I/O 部分就是解决 CPU 快、设备慢、设备复杂且共享的问题。DMA/通道减少 CPU 搬运，缓冲平滑速度差，SPOOLing 把独占设备虚拟成共享设备。}
\exam{常考设备分类、四种 I/O 控制方式、DMA 与中断/通道区别、缓冲和 SPOOLing 的作用。}


\concept{通道控制方式}
\begin{points}
  \item CPU 把一组 I/O 操作写成通道程序交给通道执行。
  \item 通道独立控制设备完成复杂 I/O 序列，完成后中断 CPU。
  \item 与 DMA 相比，通道更像专用 I/O 处理机，控制能力更强。
\end{points}
\plain{I/O 部分就是解决 CPU 快、设备慢、设备复杂且共享的问题。DMA/通道减少 CPU 搬运，缓冲平滑速度差，SPOOLing 把独占设备虚拟成共享设备。}
\exam{常考设备分类、四种 I/O 控制方式、DMA 与中断/通道区别、缓冲和 SPOOLing 的作用。}


\section{I/O 软件原理}
\concept{I/O 软件层次}
\begin{points}
  \item 用户层 I/O 软件：库函数、假脱机程序等，为用户提供高级接口。
  \item 设备独立性软件：提供统一逻辑设备接口、命名、保护、缓冲、错误处理。
  \item 设备驱动程序：针对具体设备控制器发命令、处理中断。
  \item 中断处理程序：响应设备中断，完成状态检查、唤醒等待进程等。
\end{points}
\plain{I/O 部分就是解决 CPU 快、设备慢、设备复杂且共享的问题。DMA/通道减少 CPU 搬运，缓冲平滑速度差，SPOOLing 把独占设备虚拟成共享设备。}
\exam{常考设备分类、四种 I/O 控制方式、DMA 与中断/通道区别、缓冲和 SPOOLing 的作用。}


\concept{设备独立性}
\begin{points}
  \item 设备独立性又称设备无关性，指用户程序不依赖具体物理设备。
  \item 第一层含义：独立于同类设备的具体设备号，例如打印到逻辑打印机而非某台具体打印机。
  \item 第二层含义：在一定抽象下独立于设备类型，例如统一用文件接口访问普通文件和部分设备文件。
  \item 好处：提高可移植性、灵活性和设备分配效率。
\end{points}
\plain{I/O 部分就是解决 CPU 快、设备慢、设备复杂且共享的问题。DMA/通道减少 CPU 搬运，缓冲平滑速度差，SPOOLing 把独占设备虚拟成共享设备。}
\exam{常考设备分类、四种 I/O 控制方式、DMA 与中断/通道区别、缓冲和 SPOOLing 的作用。}


\section{缓冲技术管理}
\concept{为什么需要缓冲}
\begin{points}
  \item CPU 和 I/O 设备速度差异巨大，缓冲可缓和二者速度不匹配。
  \item 缓冲可减少中断次数和设备启动次数，提高并行程度。
  \item 缓冲还能适配不同数据粒度，例如用户按字节读写，磁盘按块传输。
\end{points}
\plain{I/O 部分就是解决 CPU 快、设备慢、设备复杂且共享的问题。DMA/通道减少 CPU 搬运，缓冲平滑速度差，SPOOLing 把独占设备虚拟成共享设备。}
\exam{常考设备分类、四种 I/O 控制方式、DMA 与中断/通道区别、缓冲和 SPOOLing 的作用。}


\concept{单缓冲}
\begin{points}
  \item 系统设置一个缓冲区，设备先把数据送入缓冲区，再由进程取走。
  \item 相比无缓冲，可让设备传输和进程处理部分重叠。
  \item 但一个缓冲区仍可能使设备或进程等待。
\end{points}
\plain{I/O 部分就是解决 CPU 快、设备慢、设备复杂且共享的问题。DMA/通道减少 CPU 搬运，缓冲平滑速度差，SPOOLing 把独占设备虚拟成共享设备。}
\exam{常考设备分类、四种 I/O 控制方式、DMA 与中断/通道区别、缓冲和 SPOOLing 的作用。}


\concept{双缓冲}
\begin{points}
  \item 设置两个缓冲区，一个用于设备填充，另一个供进程处理。
  \item 可以更好地实现输入/输出与处理并行。
  \item 常用于数据流连续、生产和消费速度接近的场景。
\end{points}
\plain{I/O 部分就是解决 CPU 快、设备慢、设备复杂且共享的问题。DMA/通道减少 CPU 搬运，缓冲平滑速度差，SPOOLing 把独占设备虚拟成共享设备。}
\exam{常考设备分类、四种 I/O 控制方式、DMA 与中断/通道区别、缓冲和 SPOOLing 的作用。}


\concept{循环缓冲与缓冲池}
\begin{points}
  \item 循环缓冲用多个缓冲区构成环形队列，适合连续数据流。
  \item 缓冲池由系统统一管理多个缓冲区，按需分配和回收。
  \item 缓冲池能提高缓冲区利用率，适合多设备、多进程共享。
\end{points}
\plain{I/O 部分就是解决 CPU 快、设备慢、设备复杂且共享的问题。DMA/通道减少 CPU 搬运，缓冲平滑速度差，SPOOLing 把独占设备虚拟成共享设备。}
\exam{常考设备分类、四种 I/O 控制方式、DMA 与中断/通道区别、缓冲和 SPOOLing 的作用。}


\concept{缓冲与缓存区别}
\begin{points}
  \item 缓冲 buffer 主要解决速度不匹配和数据传输粒度不一致。
  \item 缓存 cache 主要利用局部性保存数据副本，减少慢速设备访问。
  \item 例子：磁盘块缓存是 cache；键盘输入缓冲区更偏 buffer。
\end{points}
\plain{I/O 部分就是解决 CPU 快、设备慢、设备复杂且共享的问题。DMA/通道减少 CPU 搬运，缓冲平滑速度差，SPOOLing 把独占设备虚拟成共享设备。}
\exam{常考设备分类、四种 I/O 控制方式、DMA 与中断/通道区别、缓冲和 SPOOLing 的作用。}


\section{设备分配与 SPOOLing}
\concept{设备分配中的数据结构}
\begin{points}
  \item 设备控制表 DCT：描述一台设备的状态、类型、队列等。
  \item 控制器控制表 COCT：描述设备控制器状态。
  \item 通道控制表 CHCT：描述通道状态。
  \item 系统设备表 SDT：记录系统中所有设备及其控制表入口。
\end{points}
\plain{I/O 部分就是解决 CPU 快、设备慢、设备复杂且共享的问题。DMA/通道减少 CPU 搬运，缓冲平滑速度差，SPOOLing 把独占设备虚拟成共享设备。}
\exam{常考设备分类、四种 I/O 控制方式、DMA 与中断/通道区别、缓冲和 SPOOLing 的作用。}


\concept{设备分配过程}
\begin{points}
  \item 根据用户请求的逻辑设备名查找系统设备表。
  \item 分配设备、控制器和通道，若某一级不可用则进程等待。
  \item 分配成功后启动 I/O 操作，I/O 完成后释放相关资源并唤醒等待进程。
\end{points}
\plain{I/O 部分就是解决 CPU 快、设备慢、设备复杂且共享的问题。DMA/通道减少 CPU 搬运，缓冲平滑速度差，SPOOLing 把独占设备虚拟成共享设备。}
\exam{常考设备分类、四种 I/O 控制方式、DMA 与中断/通道区别、缓冲和 SPOOLing 的作用。}


\concept{SPOOLing 系统}
\begin{points}
  \item SPOOLing 是外部设备联机并行操作技术，常用于把独占设备改造成共享设备。
  \item 基本思想：把慢速独占设备的输入/输出先放到磁盘上的输入井/输出井中，由后台进程统一和实际设备交互。
  \item 打印机例子：用户进程把打印数据写入磁盘输出井后即可继续执行，打印进程再按队列送往物理打印机。
  \item SPOOLing 组成：输入井、输出井、输入进程、输出进程、井管理程序。
  \item 直觉：用户不是直接占用打印机，而是把任务交到“打印队列”；物理打印机仍独占，但逻辑上多个用户可同时提交打印任务。
\end{points}
\plain{I/O 部分就是解决 CPU 快、设备慢、设备复杂且共享的问题。DMA/通道减少 CPU 搬运，缓冲平滑速度差，SPOOLing 把独占设备虚拟成共享设备。}
\exam{常考设备分类、四种 I/O 控制方式、DMA 与中断/通道区别、缓冲和 SPOOLing 的作用。}


\chapter*{总复习：高频辨析与做题模板}
\addcontentsline{toc}{chapter}{总复习：高频辨析与做题模板}

\section*{一、概念定位模板}
\addcontentsline{toc}{section}{一、概念定位模板}
\begin{points}
  \item 问 CPU 给谁用：进程调度。
  \item 问多个进程怎么不抢坏共享变量：同步互斥。
  \item 问互相等资源卡住：死锁。
  \item 问地址怎么变、页表、碎片：存储器管理。
  \item 问缺页、页面置换、抖动：虚拟存储器。
  \item 问文件名、目录、逻辑结构、访问方法：文件系统接口。
  \item 问文件块放磁盘哪里、inode、空闲块：文件系统实现。
  \item 问磁盘臂怎么移动、访问时间：大容量存储。
  \item 问 DMA、缓冲、SPOOLing、设备分配：设备管理。
\end{points}

\section*{二、地址转换题模板}
\addcontentsline{toc}{section}{二、地址转换题模板}
\begin{points}
  \item 看页大小/段大小，确定偏移位数或偏移范围。
  \item 把逻辑地址拆成页号/段号和偏移。
  \item 查页表/段表，判断是否越界或缺页。
  \item 拼物理地址：页式为物理块号 + 页内偏移；段式为段基址 + 段内偏移。
\end{points}

\section*{三、PV 题模板}
\addcontentsline{toc}{section}{三、PV 题模板}
\begin{points}
  \item 找共享资源：需要 mutex 互斥。
  \item 找先后关系：先发生的一方 V，后发生的一方 P。
  \item 找数量约束：空缓冲区用 empty，满缓冲区用 full，资源个数就是信号量初值。
  \item 写顺序：先等条件，再进互斥；先出互斥，再释放条件。
  \item 检查死锁：有没有拿着 mutex 等待别人改变条件。
\end{points}

\section*{四、调度题模板}
\addcontentsline{toc}{section}{四、调度题模板}
\begin{points}
  \item 列表：进程、到达时间、服务时间、优先级/队列。
  \item 找调度点：到达、完成、时间片到、阻塞、抢占、优先级调整。
  \item 画甘特图。
  \item 算完成时间、周转时间、等待时间、带权周转时间。
  \item 抢占式算法每个新进程到达都要判断是否抢占。
\end{points}

\section*{五、页面置换题模板}
\addcontentsline{toc}{section}{五、页面置换题模板}
\begin{points}
  \item 初始化页框为空，按引用串逐个访问。
  \item 命中：不计缺页，但更新 LRU/Clock 等状态。
  \item 缺页：若有空页框直接装入；否则按算法选牺牲页。
  \item 记录每一步页框内容和缺页次数。
  \item 判断 Belady：只针对 FIFO 这类算法可能出现，不要泛化到 LRU/OPT。
\end{points}

\section*{六、文件系统计算题模板}
\addcontentsline{toc}{section}{六、文件系统计算题模板}
\begin{points}
  \item 目录检索题：先算每块可放多少目录项，再算平均需要扫描多少块。
  \item 索引最大文件题：先算每个索引块可放多少块号，再根据索引级数求可寻址数据块数。
  \item 链式分配随机访问题：访问第 $k$ 块通常要从头顺链访问到第 $k$ 块。
  \item 连续分配题：直接定位强，但增长和外部碎片是弱点。
\end{points}

\section*{七、最后一遍背诵关键词}
\addcontentsline{toc}{section}{七、最后一遍背诵关键词}
\begin{multicols}{2}
\begin{points}
  \item OS：资源管理器、控制程序、接口。
  \item 四特征：并发、共享、虚拟、异步。
  \item 接口：命令接口、程序接口、图形接口。
  \item 进程：程序一次执行，PCB 是唯一标志。
  \item 线程：CPU 调度基本单位，共享进程资源。
  \item 调度：高级进系统，中级进内存，低级上 CPU。
  \item 临界区：空闲让进、忙则等待、有限等待、让权等待。
  \item 信号量：P 先减，V 先加。
  \item 死锁：互斥、请求保持、不可剥夺、循环等待。
  \item 分页：逻辑连续、物理离散、页表映射。
  \item 分段：按逻辑结构，便于共享保护。
  \item 虚存：局部性、请求调入、页面置换。
  \item 抖动：缺页太多，系统忙于换页。
  \item 文件：命名信息项序列，目录用于检索。
  \item inode：文件元数据和块指针。
  \item 磁盘：寻道 + 旋转延迟 + 传输。
  \item I/O：DMA、通道、缓冲、SPOOLing。
\end{points}
\end{multicols}

\end{document}
